\chapter{Элементы теории чисел и шифрование}
\label{ch:numtheory}

«Математика~--- царица наук, а теория чисел~--- царица математики»,~---
говорил Карл Фридрих Гаусс почти двести лет назад. Долгое время теория
чисел считалась самой «чистой» из математических дисциплин~--- занимающейся
ради красоты и никак не связанной с практикой. Сегодня, в эпоху Интернета,
банковских карт и мессенджеров, эта наука оказалась едва ли не самой
востребованной: без неё не отправить безопасное сообщение, не оплатить
покупку и не зайти на сайт по защищённому соединению.

В этой главе мы пройдём путь от античного алгоритма Евклида до современных
криптографических протоколов, обеспечивающих безопасность в Интернете.
Один и тот же небольшой набор фактов о \emph{делимости и сравнениях по
модулю} окажется работающим везде~--- от шифрования RSA до структуры
данных для подсчёта различных элементов в больших потоках.

Читатель встретит три «слоя» материала:
\begin{itemize}\itemsep1pt
  \item \textbf{Теоретико-числовой фундамент}: алгоритм Евклида,
        малая теорема Ферма, функция Эйлера, тесты простоты~---
        классические результаты, без которых не существует современной
        криптографии.
  \item \textbf{Криптографические протоколы}: симметричное и
        асимметричное шифрование, RSA и Диффи--Хеллман, электронная
        цифровая подпись, схемы электронного голосования.
  \item \textbf{Вероятностные алгоритмы}: универсальное хеширование,
        счётчики HyperLogLog, измерение «социального расстояния»
        в графах больших сетей.
\end{itemize}

Многие из изученных конструкций ежесекундно работают внутри устройств
читателя: каждое HTTPS-соединение, каждый клик в банковском приложении,
каждый запрос в поисковике~--- результат применения математических
фактов, которым посвящена эта глава.

% =============================================================
%  § 1.  Элементы теории чисел
% =============================================================

\section{Элементы теории чисел}
\label{sec:numtheory}

\subsection{Делимость, остатки и сравнения по модулю}

Целые числа --- самые знакомые объекты в математике, и кажется, что про них трудно сказать что-то новое. Однако именно из вопросов о \emph{делимости} целых чисел вырос целый раздел математики --- \term{теория чисел}, --- и именно её результаты сейчас работают внутри каждого мессенджера и банковской карты.

\begin{definition}{}{def:p1-1}
Говорят, что целое число $a$ \term{делится} на целое число $b \neq 0$ (или что $b$ \term{делит} $a$), если существует такое целое число $q$, что $a = b\cdot q$. Пишут $b \mid a$.
\end{definition}

Если $a$ не делится на $b$, то всё равно можно записать
\[
a = b\cdot q + r, \qquad 0 \le r < |b|.
\]
Это знакомое со школы \term{деление с остатком}. Число $q$ называется \emph{неполным частным}, а $r$ --- \term{остатком от деления} $a$ на $b$. Остаток обозначают также $a \bmod b$ (читается «$a$ по модулю $b$»).

\begin{example}{}{ex:p1-1}
$17 = 5\cdot 3 + 2$, поэтому $17 \bmod 5 = 2$.\\
$-17 = 5\cdot(-4) + 3$, поэтому $(-17)\bmod 5 = 3$ (остаток --- неотрицательное число!).
\end{example}

Простые числа --- «атомы» делимости.

\begin{definition}{}{def:p1-2}
Натуральное число $p > 1$ называется \term{простым}, если оно делится только на единицу и на само себя. Натуральное $n > 1$, не являющееся простым, называется \term{составным}.
\end{definition}

Первые простые числа: $2, 3, 5, 7, 11, 13, 17, 19, 23, 29, \ldots$. Простых чисел бесконечно много --- это доказал ещё Евклид (доказательство мы повторим в п.~\ref{subsec:euclid}).

Основная теорема арифметики говорит, что любое натуральное число $n > 1$ \emph{единственным} (с точностью до порядка) способом раскладывается в произведение простых:
\[
n = p_1^{\alpha_1} \cdot p_2^{\alpha_2} \cdots p_k^{\alpha_k}.
\]

\paragraph*{Сравнения по модулю.} В каждой задаче, где важна только «остаточная» информация (например, день недели в зависимости от номера дня в году), удобно вместо чисел рассматривать их остатки от деления на некоторое фиксированное число.

\begin{definition}{}{def:p1-3}
Пусть $m \ge 2$ --- натуральное число (\term{модуль}). Целые числа $a$ и $b$ называются \term{сравнимыми по модулю $m$}, если их разность $a - b$ делится на $m$. Пишут
\[
a \equiv b \modn{m}.
\]
\end{definition}

\noindent
Равносильное определение: $a \equiv b \modn{m}$ тогда и только тогда, когда $a$ и $b$ дают одинаковые остатки при делении на~$m$.

\begin{example}{}{ex:p1-2}
$23 \equiv 9 \modn 7$, потому что $23 - 9 = 14 = 2\cdot 7$.\\ В обоих случаях остаток от деления на 7 равен 2.
\end{example}

Сравнения --- очень удобный язык, потому что они \emph{ведут себя как обычные равенства}: их можно складывать, вычитать, умножать и возводить в степень.

\begin{theorem}{свойства сравнений}{thm:p1-1}
Если $a \equiv b \modn m$ и $c \equiv d \modn m$, то:
\begin{enumerate}
  \item $a + c \equiv b + d \modn m$;
  \item $a - c \equiv b - d \modn m$;
  \item $a \cdot c \equiv b \cdot d \modn m$;
  \item $a^n \equiv b^n \modn m$ для любого натурального $n$.
\end{enumerate}
\end{theorem}

\begin{proof}[Идея доказательства]
По определению $a - b = m\cdot k_1$ и $c - d = m\cdot k_2$. Тогда $(a + c) - (b + d) = m(k_1 + k_2)$ делится на $m$. Для умножения: $ac - bd = ac - bc + bc - bd = c(a-b) + b(c-d)$, что тоже делится на $m$. Возведение в степень --- следствие умножения. \qedhere
\end{proof}

С \emph{делением} ситуация сложнее. Нельзя «просто так» поделить обе части сравнения --- например, $6 \equiv 0 \modn 6$, но «сократив на 2», получим $3 \equiv 0 \modn 6$, что неверно. Когда деление законно --- мы разберём чуть позже, после знакомства с алгоритмом Евклида.

\begin{interestbox}
Сравнения встречаются повсюду в жизни. Часы --- арифметика по модулю 12 (или 24): если сейчас 10 часов, то через 5 часов будет $10 + 5 = 15 \equiv 3 \modn{12}$. День недели --- арифметика по модулю 7. Контрольная цифра ИНН, номер банковской карты, штрихкод книги --- все они вычисляются как сумма цифр (с разными весами) по некоторому модулю.
\end{interestbox}

% -----------------------------------------------------------------
\subsection{Алгоритм Евклида}
\label{subsec:euclid}

Одна из главных операций в теории чисел --- нахождение \term{наибольшего общего делителя} (НОД) двух чисел.

\begin{definition}{}{def:p1-4}
Наибольший общий делитель чисел $a$ и $b$ (не равных одновременно нулю) --- это наибольшее натуральное число, на которое делятся одновременно и $a$, и $b$. Обозначается $\NOD(a,b)$.

Если $\NOD(a,b) = 1$, то числа $a$ и $b$ называются \term{взаимно простыми}.
\end{definition}

Как найти $\NOD(a,b)$? Можно, например, разложить оба числа на простые множители и взять общие. Но при больших числах разложение на множители --- очень тяжёлая задача (об этом мы поговорим в конце параграфа). К счастью, ещё в III веке до н. э. Евклид предложил способ, который работает несравнимо быстрее.

\paragraph*{Геометрическая идея.} Представим себе прямоугольник со сторонами $a$ и $b$, причём $a > b$. Будем «выкладывать» в него квадраты со стороной $b$, пока это возможно. Когда останется полоска со сторонами $b$ и $r = a \bmod b$, повторим процедуру для неё, и так далее. Когда выкладывание закончится точным квадратом --- сторона этого последнего квадрата и есть $\NOD(a,b)$ (рис.~\ref{fig:euclid-geom}).

\begin{figure}[H]
\centering
\begin{tikzpicture}[scale=0.55, every node/.style={font=\small}]
  % Параметры для прямоугольника 48 x 30
  % 48 = 1*30 + 18 (квадрат 30x30, остаток 30x18)
  % 30 = 1*18 + 12 (квадрат 18x18, остаток 18x12)
  % 18 = 1*12 + 6  (квадрат 12x12, остаток 12x6)
  % 12 = 2*6 + 0   (два квадрата 6x6) -> НОД = 6
  \def\W{16} % 48/3
  \def\H{10} % 30/3

  % Основной прямоугольник 48x30 -> 16x10
  \draw[thick, prosvBlue] (0,0) rectangle (\W,\H);

  % Квадрат 30x30 -> 10x10 (слева)
  \fill[prosvLight, opacity=0.7] (0,0) rectangle (10,10);
  \draw[thick, prosvBlue] (10,0) -- (10,10);
  \node at (5,5) {\large $30\times 30$};

  % Остаток 18x30 -> 6x10 (справа). Делим: квадрат 18x18 -> 6x6 сверху
  \fill[prosvCream, opacity=0.8] (10,4) rectangle (16,10);
  \draw[thick, prosvBlue] (10,4) -- (16,4);
  \node at (13,7) {$18\times 18$};

  % Остаток 18x12 -> 6x4 снизу. Квадрат 12x12 -> 4x4 слева
  \fill[prosvLight!60!white, opacity=0.7] (10,0) rectangle (14,4);
  \draw[thick, prosvBlue] (14,0) -- (14,4);
  \node at (12,2) {$12\!\times\!12$};

  % Остаток 6x4 -> 2x4 справа. Делим: два квадрата 6x6 = 2x2
  \fill[prosvGold!30!white, opacity=0.8] (14,0) rectangle (16,2);
  \draw[thick, prosvBlue] (14,2) -- (16,2);
  \draw[thick, prosvBlue] (15,0) -- (15,2);
  \fill[prosvGold!30!white, opacity=0.8] (14,2) rectangle (16,4);
  \draw[thick, prosvBlue] (15,2) -- (15,4); 

  \node[font=\scriptsize] at (15,1) {6};
  \node[font=\scriptsize] at (15,3) {6};
  
  % Подписи сторон
  \draw[<->,>=stealth, prosvRed] (0,-0.6) -- (\W,-0.6);
  \node[below, prosvRed] at (\W/2,-0.7) {$a = 48$};
  
  \draw[<->,>=stealth, prosvRed] (-0.6,0) -- (-0.6,\H);
  \node[left, prosvRed] at (-0.7,\H/2) {$b = 30$};

  \draw[<->,>=stealth, prosvGreen] (16.4,0) -- (16.4,2);
  \node[right, prosvGreen] at (16.5,1) {$\NOD = 6$};
\end{tikzpicture}
\caption{Геометрическая иллюстрация алгоритма Евклида для пары $(48,\,30)$. Прямоугольник заполняется квадратами: $30{\times}30$, затем $18{\times}18$, затем $12{\times}12$ и, наконец, два квадрата $6{\times}6$. Сторона последнего квадрата равна $\NOD(48,30) = 6$.}
\label{fig:euclid-geom}
\end{figure}

Численно это соответствует последовательным делениям с остатком:
\begin{align*}
48 &= 1\cdot 30 + 18,\\
30 &= 1\cdot 18 + 12,\\
18 &= 1\cdot 12 + 6,\\
12 &= 2\cdot 6 + 0 \;\;\Longrightarrow\;\; \NOD(48,30) = 6.
\end{align*}

\noindent
\emph{Почему это работает?} Ключевое наблюдение: если $a = bq + r$, то $\NOD(a, b) = \NOD(b, r)$. Действительно, любое число, делящее $a$ и $b$, делит и $r = a - bq$; обратно, любое число, делящее $b$ и $r$, делит и $a = bq + r$. Значит, у пар $(a,b)$ и $(b,r)$ один и тот же набор общих делителей --- и, следовательно, один и тот же $\NOD$. На каждом шаге второй аргумент уменьшается; раньше или позже он станет нулём, и тогда первое число --- искомый НОД.

\begin{algorithmbox}[title={Алгоритм: нахождение НОД (Евклид)}]
\textbf{Вход:} натуральные числа $a, b$, $a \ge b > 0$.\\
\textbf{Выход:} $\NOD(a,b)$.
\begin{enumerate}
  \item Если $b = 0$, вернуть $a$.
  \item Иначе вычислить $r = a \bmod b$ и повторить с парой $(b, r)$.
\end{enumerate}
\end{algorithmbox}

На псевдокоде:
\begin{center}
\begin{minipage}{0.7\linewidth}
\ttfamily\small
\noindent
function gcd(a, b):\\
\quad while b $\ne$ 0:\\
\qquad a, b := b, a mod b\\
\quad return a
\end{minipage}
\end{center}

\paragraph*{Скорость работы.} Сколько шагов делает алгоритм? Можно доказать (теорема Ламе, 1844 г.): число делений в алгоритме Евклида для $\NOD(a,b)$ не превосходит примерно $5\log_{10}\min(a,b)$. Это значит, что для чисел в \emph{сотни цифр} (а такие числа реально используются в криптографии) НОД находится за несколько сотен делений --- доли секунды.

% -----------------------------------------------------------------
\subsection{Расширенный алгоритм Евклида}

Поразительный факт: Евклид нашёл не только $\NOD$, но и способ \emph{представить} его в виде «целочисленной линейной комбинации» исходных чисел.

\begin{theorem}{соотношение Безу}{thm:p1-2}
\label{thm:bezout}
Для любых натуральных $a, b$ существуют целые числа $x, y$ (\emph{не обязательно положительные}), такие что
\[
a\cdot x + b\cdot y = \NOD(a,b).
\]
\end{theorem}

Эти $x, y$ называются \term{коэффициентами Безу}. Найти их позволяет \term{расширенный алгоритм Евклида}: при делениях с остатком мы заодно «отслеживаем», как выражается каждый текущий остаток через исходные $a$ и $b$.

\paragraph*{Пример: те же 48 и 30.} Прокрутим цепочку делений «снизу вверх», выражая 6 (наш НОД) через всё бо\'{}льшие числа.
\begin{align*}
6 &= 18 - 1\cdot 12 && \text{(из третьего деления)}\\
  &= 18 - 1\cdot (30 - 1\cdot 18) = 2\cdot 18 - 1\cdot 30 && \text{(подставили } 12 = 30 - 18 \text{)}\\
  &= 2\cdot (48 - 1\cdot 30) - 1\cdot 30 = 2\cdot 48 - 3\cdot 30 && \text{(подставили } 18 = 48 - 30 \text{)}
\end{align*}
Получили: $\boxed{2\cdot 48 + (-3)\cdot 30 = 6}$. Коэффициенты Безу: $x = 2$, $y = -3$.

\begin{example}{}{ex:p1-3}
Проверим: $2\cdot 48 - 3\cdot 30 = 96 - 90 = 6 = \NOD(48,30)$. Совпадает.
\end{example}

Удобно вести расчёт в таблице:
\begin{center}
\begin{tabular}{|c|c|c|c|c|}
\hline
шаг & $r$ & $q$ & $x$ & $y$ \\\hline
0 & 48 & --- & 1 & 0\\
1 & 30 & --- & 0 & 1\\
2 & 18 & 1 & 1 & $-1$\\
3 & 12 & 1 & $-1$ & 2\\
4 & 6  & 1 & 2 & $-3$\\
5 & 0  & 2 & --- & ---\\\hline
\end{tabular}
\end{center}
В каждой строке выполняются равенства $r = 48 x + 30 y$. На предпоследней строке (где $r = \NOD = 6$) и считываются коэффициенты.

% -----------------------------------------------------------------
\subsection{Решение линейных сравнений}

Применим то, что получили, к решению уравнений в модулярной арифметике.

\paragraph*{Обратный элемент по модулю.} В обычной арифметике число, обратное к $a$, --- это $1/a$: при умножении даёт единицу. По модулю $m$ \term{обратным к $a$} называется такое целое $a^{-1}$, что
\[
a \cdot a^{-1} \equiv 1 \modn m.
\]
Например, по модулю 7: $3\cdot 5 = 15 = 2\cdot 7 + 1 \equiv 1 \modn 7$, значит, $3^{-1} \equiv 5 \modn 7$.

\begin{theorem}{}{thm:p1-5}
Обратный элемент к $a$ по модулю $m$ \emph{существует} тогда и только тогда, когда $\NOD(a, m) = 1$. При этом он находится с помощью расширенного алгоритма Евклида.
\end{theorem}

\begin{proof}[Идея]
Если $\NOD(a,m) = 1$, то по теореме~\ref{thm:bezout} существуют целые $x, y$, такие что $ax + my = 1$. Это сравнение по модулю $m$ принимает вид $ax \equiv 1 \modn m$, то есть $x \equiv a^{-1} \modn m$.

Обратно, если $a\cdot a^{-1} \equiv 1 \modn m$, то $a \cdot a^{-1} - 1 = m\cdot k$ для некоторого $k$, откуда любой общий делитель $a$ и $m$ должен делить $1$, --- значит, $\NOD(a,m) = 1$.
\end{proof}

\begin{example}{}{ex:p1-4}
Найдём $11^{-1} \modn{26}$. Применяем расширенный алгоритм Евклида к $(26,11)$:
\begin{align*}
26 &= 2\cdot 11 + 4,\\
11 &= 2\cdot 4 + 3,\\
4 &= 1\cdot 3 + 1,\\
3 &= 3\cdot 1 + 0.
\end{align*}
$\NOD = 1$. Раскручиваем:
\begin{align*}
1 &= 4 - 1\cdot 3 = 4 - (11 - 2\cdot 4) = 3\cdot 4 - 11\\
&= 3\cdot (26 - 2\cdot 11) - 11 = 3\cdot 26 - 7\cdot 11.
\end{align*}
Значит, $-7\cdot 11 \equiv 1 \modn{26}$, то есть $11^{-1} \equiv -7 \equiv 19 \modn{26}$. Проверка: $11\cdot 19 = 209 = 8\cdot 26 + 1$.
\end{example}

\paragraph*{Линейные сравнения.} Уравнение вида
\[
a\cdot x \equiv b \modn m
\]
называется \term{линейным сравнением}. Если $\NOD(a,m) = 1$, его решение находится «домножением на обратный»: $x \equiv a^{-1}\cdot b \modn m$.

% -----------------------------------------------------------------
\subsection{Быстрое возведение в степень}
\label{subsec:fastpow}

В криптографии нужно уметь вычислять выражения вида $a^n \bmod m$ для огромных $n$ (порядка $2^{1000}$). Если в лоб умножать $a$ на себя $n$ раз --- это никогда не закончится: даже за миллиард умножений в секунду на это ушло бы больше времени, чем существует Вселенная.

Спасает идея \term{«разделяй и властвуй»}, известная также как \term{бинарное возведение в степень}. Идея проста: разделим показатель пополам:
\[
a^n = \begin{cases}
\bigl(a^{n/2}\bigr)^2, & n\text{ чётно},\\[2pt]
a\cdot \bigl(a^{(n-1)/2}\bigr)^2, & n\text{ нечётно}.
\end{cases}
\]
Каждое такое деление пополам уменьшает показатель примерно вдвое, поэтому всего нужно около $\log_2 n$ шагов вместо~$n$ (рис.~\ref{fig:binpow}).

\begin{figure}[H]
\centering
\begin{tikzpicture}[
  level distance=1.0cm,
  level 1/.style={sibling distance=10cm},
  level 2/.style={sibling distance=5cm},
  level 3/.style={sibling distance=2.5cm},
  every node/.style={
    rounded corners=2pt,
    draw=prosvBlue,
    fill=prosvLight,
    font=\small,
    inner sep=3pt
  },
  edge from parent/.style={->, >=stealth, draw=prosvBlue, thick}
]
\node {$a^{13}$}
  child {node {$a^{13} = a\cdot a^{12}$}
    child {node {$a^{12} = (a^6)^2$}
      child {node {$a^6 = (a^3)^2$}
        child {node {$a^3 = a\cdot a^2$}}
      }
    }
  };
\end{tikzpicture}
\caption{Дерево бинарного возведения $a^{13}$: показатель уменьшается вдвое (с поправкой на чётность) на каждом шаге. Всего около $\log_2 13 \approx 4$ шагов.}
\label{fig:binpow}
\end{figure}

\noindent
Каждый «спуск» делит показатель примерно на 2, поэтому всего получается около $\log_2 n$ шагов. Для $n = 2^{1000}$ это \emph{тысяча} умножений вместо $2^{1000}$.

На псевдокоде:
\begin{center}
\begin{minipage}{0.78\linewidth}
\ttfamily\small
\noindent
function powmod(a, n, m):\\
\quad result := 1\\
\quad a := a mod m\\
\quad while n > 0:\\
\qquad if n is odd:\\
\quad\qquad result := (result $\ast$ a) mod m\\
\qquad a := (a $\ast$ a) mod m\\
\qquad n := n div 2\\
\quad return result
\end{minipage}
\end{center}

\begin{example}{}{ex:p1-5}
Вычислим $7^{13} \bmod 23$. Запишем $13 = 1101_2$, то есть $13 = 8 + 4 + 1$. Значит, $7^{13} = 7^{8}\cdot 7^{4}\cdot 7^{1}$. Считаем последовательно по модулю 23:
\[
7^1 = 7,\quad 7^2 = 49 \equiv 3,\quad 7^4 \equiv 3^2 = 9,\quad 7^8 \equiv 9^2 = 81 \equiv 81 - 3\cdot 23 = 12.
\]
Тогда $7^{13} \equiv 12\cdot 9\cdot 7 = 756 \equiv 756 - 32\cdot 23 = 756 - 736 = 20 \modn{23}$. Всего --- около пяти умножений вместо тринадцати.
\end{example}

\paragraph*{Та же идея в другом месте: числа Фибоначчи.} Знакомая со школы последовательность Фибоначчи задаётся рекуррентно:
\[
F_0 = 0,\quad F_1 = 1,\quad F_{n+1} = F_n + F_{n-1}.
\]

На первый взгляд, чтобы вычислить $F_n$, нужно сделать $n$ сложений --- линейное время. Однако и здесь работает «разделяй и властвуй»! Запишем рекуррентное соотношение в матричной форме:
\[
\begin{pmatrix} F_{n+1} \\ F_n \end{pmatrix}
\;=\;
\begin{pmatrix} 1 & 1 \\ 1 & 0 \end{pmatrix}
\begin{pmatrix} F_n \\ F_{n-1} \end{pmatrix}.
\]
Обозначим матрицу как $A$. Применяя её $n$ раз к вектору $(F_1, F_0)^\top = (1,0)^\top$, получим:
\[
\begin{pmatrix} F_{n+1} \\ F_n \end{pmatrix}
\;=\;
A^n
\begin{pmatrix} 1 \\ 0 \end{pmatrix}.
\]
Прямое доказательство --- индукцией по $n$.

\begin{importantbox}
Тем самым задача сводится к возведению матрицы $2{\times}2$ в степень $n$. Но матрицы можно перемножать «делением пополам» точно так же, как числа! Сложность --- $O(\log n)$ матричных умножений, то есть около $\log_2 n$ операций; вместо линейного числа шагов получается логарифмическое (рис.~\ref{fig:p1-3}).
\end{importantbox}

\begin{figure}[H]
\centering
\begin{tikzpicture}[scale=1.0, every node/.style={font=\small}]
  % Сравнение линейного и логарифмического подхода
  \begin{scope}[xshift=-5cm]
    \node[font=\bfseries\color{prosvRed}] at (1.8,2.4) {Наивно: $n$ шагов};
    \foreach \i in {0,...,9}{
      \fill[prosvRed!60!white] (\i*0.35,0) rectangle (\i*0.35+0.3,0.5);
    }
    \node at (1.5, -0.4) {$F_1 \to F_2 \to \cdots \to F_{10}$};
  \end{scope}
  \begin{scope}[xshift=1cm]
    \node[font=\bfseries\color{prosvGreen}] at (1.8,2.4) {С «делением пополам»: $\log_2 n$ шагов};
    \foreach \i in {0,...,3}{
      \fill[prosvGreen!70!white] (\i*0.5,0) rectangle (\i*0.5+0.45,0.5);
    }
    \node at (1.5,-0.4) {всего $\approx \log_2 10 \approx 4$ шага};
  \end{scope}
\end{tikzpicture}
\caption{Сравнение наивного и «разделяй и властвуй»-подхода к вычислению $F_n$.}
\label{fig:p1-3}
\end{figure}

Этот пример замечателен тем, что показывает: даже линейная на первый взгляд задача может «прятать» внутри себя структуру, ускоряющую решение экспоненциально.

% -----------------------------------------------------------------
\subsection{Функция Эйлера и малая теорема Ферма}

Чем дальше, тем интересней. Сейчас мы введём одну функцию, играющую ключевую роль в современной криптографии.

\begin{definition}{}{def:p1-5}
\term{Функцией Эйлера} $\Eul(n)$ называется количество натуральных чисел в отрезке $[1, n]$, взаимно простых с $n$.
\end{definition}

\begin{example}{}{ex:p1-6}
$\Eul(1) = 1$, $\Eul(2) = 1$, $\Eul(6) = 2$ (а именно, числа 1 и 5), $\Eul(10) = 4$ (числа 1, 3, 7, 9).
\end{example}

\paragraph*{Случай простого числа.} Если $p$ --- простое число, то с $p$ \emph{не} взаимно прост лишь сам $p$ среди чисел $1, 2, \ldots, p$. Значит,
\[
\boxed{\;\Eul(p) = p - 1\;}\quad \text{для простого } p.
\]

\paragraph*{Случай произведения двух простых.} Пусть $n = p\cdot q$, где $p, q$ --- различные простые. Какие числа из $[1, n]$ не взаимно просты с $n$? Только кратные $p$ или кратные $q$. Кратных $p$ ровно $q$ штук (это $p, 2p, \ldots, qp$); кратных $q$ ровно $p$ штук; единственное число, кратное обоим, --- $pq = n$. По формуле включений-исключений:
\[
\Eul(pq) = pq - p - q + 1 = (p-1)(q-1).
\]
Это короткое наблюдение --- сердце алгоритма RSA, который мы изучим в §~\ref{sec:rsa-dh}.

\begin{example}{}{ex:p1-7}
$\Eul(15) = \Eul(3\cdot 5) = 2\cdot 4 = 8$. Действительно, взаимно простыми с 15 в отрезке $[1,15]$ являются числа $\{1, 2, 4, 7, 8, 11, 13, 14\}$ --- ровно 8 чисел.
\end{example}

В общем виде: для $n = p_1^{\alpha_1}\cdots p_k^{\alpha_k}$ выполняется $\Eul(n) = n\prod_{i=1}^k\!\bigl(1 - \tfrac{1}{p_i}\bigr)$, но эта формула нам сегодня понадобится только в указанных двух частных случаях.

\paragraph*{Малая теорема Ферма.} Один из самых красивых результатов классической теории чисел.

\begin{theorem}{малая теорема Ферма, 1640 г.}{thm:p1-3}
\label{thm:fermat}
Пусть $p$ --- простое число и $a$ --- целое, не делящееся на $p$. Тогда
\[
a^{p-1} \equiv 1 \modn p.
\]
Равносильно: для любого целого $a$ выполняется $a^p \equiv a \modn p$.
\end{theorem}

\begin{proof}[Схема доказательства]
Рассмотрим набор $\{a, 2a, 3a, \ldots, (p-1)a\}$ --- $(p-1)$ ненулевых остатков, умноженных на $a$. Все они различны по модулю $p$: если $ia \equiv ja \modn p$, то $a(i-j) \equiv 0 \modn p$, а так как $a$ не делится на $p$ и $|i-j| < p$, получаем $i = j$.

Значит, набор $\{a, 2a, \ldots, (p-1)a\}$ --- это в точности (с точностью до порядка) набор $\{1, 2, \ldots, p-1\}$ по модулю $p$. Перемножим всё:
\[
a\cdot 2a\cdot 3a \cdots (p-1)a \;\equiv\; 1\cdot 2\cdot 3\cdots(p-1) \modn p,
\]
то есть $a^{p-1}\cdot (p-1)! \equiv (p-1)! \modn p$. Поскольку $(p-1)!$ взаимно прост с $p$, его можно «сократить» --- получим $a^{p-1} \equiv 1 \modn p$. \qedhere
\end{proof}

Обобщение --- \term{теорема Эйлера}: $a^{\Eul(n)} \equiv 1 \modn n$ для любого $a$, взаимно простого с $n$. Доказывается аналогично, только вместо чисел $1, \ldots, p-1$ берутся все числа из $[1,n]$, взаимно простые с $n$.

\begin{interestbox}
Пьер Ферма, по профессии юрист, посвящал математике вечера. В 1640 году он написал письмо своему другу с этим утверждением, заметив: «Я бы прислал доказательство, если бы не боялся, что оно слишком длинное». Первое известное полное доказательство опубликовал Леонард Эйлер --- спустя сто лет.
\end{interestbox}

% -----------------------------------------------------------------
\subsection{Сколько вокруг простых чисел?}

Прежде чем переходить к криптографии, ответим на важный практический вопрос: \emph{часто ли встречаются простые числа?} Если простых чисел мало --- ими нельзя будет полноценно пользоваться.

\begin{theorem}{Чебышёв -- Адамар -- Валле-Пуссен, асимптотический закон распределения простых чисел}{thm:p1-4}
Пусть $\pi(N)$ --- количество простых чисел, не превосходящих $N$. Тогда
\[
\pi(N) \;\sim\; \frac{N}{\ln N} \qquad \text{при } N \to \infty.
\]
\end{theorem}

\noindent
В пересчёте на «вероятность»: если случайно взять натуральное число вокруг $N$, оно окажется простым примерно с шансом $1/\ln N$. Например, среди 1000-значных чисел простыми оказывается примерно одно из $\ln 10^{1000} = 1000\ln 10 \approx 2300$. Это совсем немало (рис.~\ref{fig:p1-4}).

\begin{figure}[H]
\centering
\begin{tikzpicture}
\begin{axis}[
  width=11cm, height=6cm,
  xlabel={$N$}, ylabel={число простых $\le N$},
  xmin=1, xmax=1000, ymin=0, ymax=200,
  axis lines=left,
  grid=major, grid style={prosvGray!30},
  legend pos=north west, legend cell align=left,
  xlabel style={font=\small}, ylabel style={font=\small},
  tick label style={font=\scriptsize},
  every axis plot/.append style={very thick}
]
% Точная функция pi(N) приближённо
\addplot[prosvRed, domain=2:1000, samples=80] {x/ln(x)};
\addlegendentry{$\dfrac{N}{\ln N}$ (приближение)};
% Фактические точки pi(N) для N=10,50,100,...
\addplot+[
  only marks, mark=*, prosvBlue, mark size=1.4pt
] coordinates {
  (10,4) (50,15) (100,25) (200,46) (300,62) (400,78) (500,95)
  (600,109) (700,125) (800,139) (900,154) (1000,168)
};
\addlegendentry{точное значение $\pi(N)$};
\end{axis}
\end{tikzpicture}
\caption{Закон распределения простых чисел: количество простых чисел до $N$ ведёт себя как $N/\ln N$.}
\label{fig:p1-4}
\end{figure}

Практический смысл: чтобы найти простое число длиной, скажем, 1024 бита для RSA-ключа, нужно перебрать в среднем порядка 700 случайных чисел и каждое проверить на простоту. И это нас приводит к следующему вопросу.

% -----------------------------------------------------------------
\subsection{Проверка простоты и задача факторизации}
\label{subsec:primality}

Два «двойственных» вопроса:
\begin{description}
  \item[Проверка простоты:] дано число $n$ --- простое оно или составное?
  \item[Факторизация:] дано составное $n$ --- разложить его на простые множители.
\end{description}

Может показаться, что эти задачи примерно одинаковой сложности. Но это не так: \emph{они принципиально разной сложности}, и именно на этом контрасте основана вся современная криптография.

\paragraph*{Тест Ферма.} Если $p$ простое, то по теореме~\ref{thm:fermat} для любого $a$, не кратного $p$, выполняется $a^{p-1} \equiv 1 \modn p$. Возьмём это как \emph{пробный признак}: если для случайного $a$ выполнено $a^{n-1} \not\equiv 1 \modn n$, то $n$ заведомо \emph{не простое}. Если же $a^{n-1} \equiv 1 \modn n$ --- то $n$ простое \emph{с большой вероятностью} (но иногда бывают и составные числа, проходящие тест --- так называемые \term{числа Кармайкла}).

\paragraph*{Тест Миллера--Рабина.} Усовершенствованный вариант теста Ферма. Он использует ту же идею, но дополнительно «отсеивает» числа Кармайкла. Принципиальные свойства:
\begin{itemize}
  \item это \term{вероятностный} тест: при отрицательном ответе число точно составное; при положительном --- простое с вероятностью ошибки не больше $1/4$ за один проход;
  \item делая $k$ независимых проходов, мы уменьшаем вероятность ошибки до $4^{-k}$. Уже при $k = 20$ это меньше, чем $10^{-12}$;
  \item работает за время $O(k\cdot \log^3 n)$ --- то есть \term{полиномиально} от числа цифр.
\end{itemize}

\paragraph*{Тест AKS (2002 г.).} Манидра Агравал, Нираж Каял и Нитин Саксена --- индийские математики --- впервые предъявили \emph{детерминированный} полиномиальный алгоритм проверки простоты. На практике он работает медленнее Миллера--Рабина, но важен как теоретическое достижение: показал, что класс задач, разрешимых за полиномиальное время, включает в себя проверку простоты.

\begin{interestbox}
Открытие AKS-алгоритма стало сенсацией: один из руководителей --- Нитин Саксена --- был на тот момент аспирантом первого года. Результат опубликован в 2004 году в одном из самых престижных математических журналов --- «Annals of Mathematics» --- и был назван «алгоритмом тысячелетия» в области теории чисел.
\end{interestbox}

\paragraph*{А что с факторизацией?} Тут картина совсем другая. Лучший известный сегодня алгоритм --- \emph{общий метод решета числового поля} --- работает за время, растущее как
\[
\exp\!\bigl(c\cdot \sqrt[3]{(\ln n)(\ln\ln n)^{2}}\bigr),
\]
что по сути \term{субэкспоненциально}, но \emph{не полиномиально}. Для 1024-битных $n$ это означает миллиарды лет работы суперкомпьютера.

\begin{importantbox}
Колоссальный разрыв в сложности: \emph{проверить, простое ли число} --- легко (полиномиально), а \emph{разложить число на простые множители} --- невероятно трудно (субэкспоненциально). Именно этот разрыв и эксплуатируется в RSA и других криптосистемах. Безопасность интернета прямо опирается на то, что у человечества пока нет быстрого алгоритма факторизации.
\end{importantbox}

\paragraph*{А если появится квантовый компьютер?} В 1994 году Питер Шор предложил квантовый алгоритм, факторизующий числа за полиномиальное время. Реальных квантовых компьютеров достаточного масштаба пока нет (на 2025 год удавалось факторизовать лишь крошечные числа в качестве демонстрации). Но если такие компьютеры будут построены --- современный RSA рухнет, и потребуются принципиально новые криптографические схемы (это направление называется \engterm{постквантовая криптография}).

% -----------------------------------------------------------------
\subsection*{Что мы получили}

Подведём итог. На вопрос «зачем нам всё это в учебнике информатики» мы готовы ответить: только что мы построили \emph{инструменты}, без которых не работает ни один защищённый сайт. А именно:
\begin{itemize}
  \item \term{сравнения по модулю} --- язык, на котором будем описывать шифрование;
  \item \term{расширенный алгоритм Евклида} --- быстро ищет обратный элемент по модулю;
  \item \term{быстрое возведение в степень} --- позволяет работать с гигантскими степенями;
  \item \term{малая теорема Ферма} и \term{функция Эйлера} --- ключ к схеме RSA;
  \item \term{разрыв в сложности} между проверкой простоты и факторизацией --- источник «безопасности».
\end{itemize}

\noindent
В следующих параграфах мы соберём из этих кирпичиков работающие криптографические протоколы.

% -----------------------------------------------------------------
\begin{tasksbox}
\begin{enumerate}
  \item Найдите $\NOD(231, 165)$ алгоритмом Евклида. Представьте НОД в виде $231 x + 165 y$.
  \item Найдите $5^{-1} \modn{37}$. Используйте расширенный алгоритм Евклида.
  \item Вычислите $3^{100} \bmod 7$, пользуясь малой теоремой Ферма.
  \item Сколько существует чисел в отрезке $[1, 100]$, взаимно простых с $100$? (Подсказка: $100 = 2^2\cdot 5^2$, но вы можете использовать формулу для $\Eul(p^2)$ или прямой подсчёт.)
  \item Запишите быстрое возведение в степень в виде программы на знакомом вам языке программирования.
  \item* Докажите, что для составного $n$ всегда найдётся $a$ из $[2, n-1]$, для которого $a^{n-1} \not\equiv 1 \modn n$ --- кроме чисел Кармайкла. Найдите наименьшее число Кармайкла.
  \item* С помощью матричной формулы для чисел Фибоначчи вычислите $F_{20}$.
\end{enumerate}
\end{tasksbox}

% =============================================================
%  § 2. Шифрование: от античности до наших дней
% =============================================================

\section{Шифрование: от тайнописи к математической задаче}
\label{sec:history}

\subsection{Зачем людям шифры}

Каждый из вас, заходя на любой сайт по адресу, начинающемуся с \texttt{https://}, доверяет ему \emph{шифрование}. Браузер и сервер обмениваются миллионами сообщений, но никто посторонний прочесть их не может. В этом параграфе мы поймём, какие задачи решает шифрование, какие общие схемы существуют и почему именно теория чисел оказалась его математической основой.

\paragraph*{Античность.} Первые шифры были чисто механическими. В Древней Греции (V век до н.\,э.) воины Спарты использовали \term{скиталу} --- цилиндр определённого диаметра, на который наматывалась лента с текстом. Без точно такой же палочки прочесть сообщение было нельзя. Юлий Цезарь, по свидетельству Светония, в переписке со своими полководцами сдвигал каждую букву алфавита на три позиции: \texttt{A → D}, \texttt{B → E}, и так далее. Этот \term{шифр Цезаря} --- простейший \emph{шифр замены} (рис.~\ref{fig:p2-1}).

\begin{figure}[H]
\centering
\begin{tikzpicture}[scale=0.85, every node/.style={font=\small\ttfamily}]
  % Открытый текст
  \foreach \i/\ch in {0/V,1/E,2/N,3/I,4/V,5/I,6/D,7/I,8/V,9/I,10/C,11/I}{
    \draw[prosvBlue] (\i*0.6, 1.2) rectangle (\i*0.6+0.55, 1.7);
    \node at (\i*0.6+0.275, 1.45) {\ch};
  }
  \node[left, font=\small\rmfamily] at (0,1.45) {Открытый текст:};
  
  % Стрелки
  \foreach \i in {0,...,11}{
    \draw[->,>=stealth, prosvGray] (\i*0.6+0.275, 1.15) -- (\i*0.6+0.275, 0.75);
  }
  
  % Шифр-текст (сдвиг +3)
  \foreach \i/\ch in {0/Y,1/H,2/Q,3/L,4/Y,5/L,6/G,7/L,8/Y,9/L,10/F,11/L}{
    \draw[prosvRed] (\i*0.6, 0.2) rectangle (\i*0.6+0.55, 0.7);
    \node at (\i*0.6+0.275, 0.45) {\ch};
  }
  \node[left, font=\small\rmfamily] at (0,0.45) {Шифр-текст:};
  
  \node[right, font=\small\rmfamily] at (7.5,1.0) {\itshape\color{prosvGray}«пришёл, увидел, победил»};
\end{tikzpicture}
\caption{Шифр Цезаря со сдвигом~3: каждая буква заменяется на букву через три позиции в алфавите.}
\label{fig:p2-1}
\end{figure}

Любой такой шифр легко взламывается: букв в алфавите немного, и можно просто перебрать все 25 возможных сдвигов. Более изобретательные шифры (например, \term{шифр Виженера} XVI века) использовали \emph{переменный} сдвиг, заданный ключевым словом. Они продержались несколько столетий, пока в середине XIX века Чарлз Бэббидж и независимо Фридрих Касиски не предложили универсальный метод их взлома.

\paragraph*{XX век: машины.} Во время Второй мировой войны на сцену вышли электромеханические шифровальные машины --- знаменитая немецкая \term{Энигма}, японская «Пурпурная» (Purple), советская «Фиалка». Их взлом стал отдельной героической страницей: в Великобритании Алан Тьюринг возглавлял проект \emph{Bletchley Park}, где была сконструирована машина «Бомба», взламывавшая Энигму. Считается, что этот успех сократил войну в Европе на 2--4 года.

\begin{interestbox}
Деятельность Алана Тьюринга и его коллег была настолько секретной, что заслуги команды Bletchley Park были рассекречены только в 1970-х годах --- то есть через 25 лет после конца войны. Сам Тьюринг при жизни так и не узнал о признании. Тьюринг также считается одним из отцов теоретической информатики --- именно его именем названа престижная Премия Тьюринга, аналог «Нобелевской премии в информатике».
\end{interestbox}

% -----------------------------------------------------------------
\subsection{Общая схема симметричного шифрования}

Все шифры до 1970-х годов имели одну общую особенность: для шифрования и расшифрования использовался \emph{один и тот же} ключ. Такие шифры называются \term{симметричными} (или \engterm{secret-key}) (рис.~\ref{fig:symenc}).

\begin{figure}[H]
\centering
\begin{tikzpicture}[
  font=\small,
  every node/.style={align=center},
  box/.style={rectangle, draw=prosvBlue, fill=prosvLight, minimum width=2.5cm, minimum height=1cm, rounded corners=2pt},
  evil/.style={rectangle, draw=prosvRed, fill=prosvCream, minimum width=2.2cm, minimum height=0.8cm, rounded corners=2pt}
]
  \node[circle, draw=prosvBlue, fill=white, minimum size=0.7cm] (a) at (-5,0) {A};
  \node[below, font=\scriptsize] at (-5,-0.5) {Алиса};
  \node[circle, draw=prosvBlue, fill=white, minimum size=0.7cm] (b) at (5,0) {B};
  \node[below, font=\scriptsize] at (5,-0.5) {Боб};
  
  \node[box] (enc) at (-2.5,0) {Шифрование $E_K$};
  \node[box] (dec) at (2.5,0) {Расшифрование $D_K$};
  
  \draw[->,>=stealth, prosvBlue, thick] (a) -- node[above]{$m$} (enc);
  \draw[->,>=stealth, prosvBlue, thick] (enc) -- node[above]{$c = E_K(m)$} (dec);
  \draw[->,>=stealth, prosvBlue, thick] (dec) -- node[above]{$m$} (b);
  
  % Ключи
  \node[font=\Large] at (-2.5, 1.2) {[K]};
  \node[font=\scriptsize] at (-2.5, 1.7) {ключ $K$};
  \node[font=\Large] at (2.5, 1.2) {[K]};
  \node[font=\scriptsize] at (2.5, 1.7) {ключ $K$};
  \draw[dashed, prosvGray] (-2.5,1.0) -- (2.5,1.0);
  
  % Ева
  \node[evil] (eve) at (0, -1.8) {Ева подслушивает\\ канал};
  \draw[->,>=stealth, prosvRed, dashed] (0, -0.3) -- (eve);
\end{tikzpicture}
\caption{Симметричное шифрование: один общий секретный ключ $K$ позволяет и зашифровать, и расшифровать. Подслушивающая Ева не должна узнать $K$.}
\label{fig:symenc}
\end{figure}

Симметричные шифры быстры и до сих пор широко применяются (например, стандарт \engterm{AES} --- то, чем шифруется содержимое любого Telegram-сообщения после установки соединения). Но у них есть фундаментальная \emph{проблема распределения ключей}: как Алисе и Бобу договориться о ключе, если они никогда раньше не встречались, а связь между ними прослушивается?

\begin{importantbox}
\textbf{Парадокс симметричного шифрования.} Чтобы безопасно обмениваться сообщениями, нужен общий ключ. Чтобы безопасно передать ключ --- нужен другой общий ключ. И так до бесконечности.
\end{importantbox}

В большом мире (миллиарды пользователей интернета) задача распределения ключей становится непреодолимой: каждой паре нужен свой ключ, всего $\binom{n}{2} \approx n^2/2$ ключей.

% -----------------------------------------------------------------
\subsection{Революция 1976 года: открытый ключ}

В 1976 году произошёл прорыв, перевернувший всю криптографию. Уитфилд Диффи и Мартин Хеллман предложили принципиально новую идею: \term{асимметричное шифрование}, или \term{криптография с открытым ключом}.

Каждый пользователь имеет \emph{пару} ключей:
\begin{itemize}
  \item \term{открытый ключ} (\engterm{public key}) --- его можно (и нужно) сообщать всем: публиковать на сайте, отправлять незашифрованно;
  \item \term{закрытый ключ} (\engterm{private key}) --- он хранится в строжайшем секрете у владельца.
\end{itemize}

Ключевое свойство: что зашифровано открытым ключом --- расшифровать может только владелец соответствующего закрытого ключа (рис.~\ref{fig:asymenc}).

\begin{figure}[H]
\centering
\begin{tikzpicture}[
  font=\small,
  every node/.style={align=center},
  box/.style={rectangle, draw=prosvBlue, fill=prosvLight, minimum width=2.4cm, minimum height=1cm, rounded corners=2pt}
]
  \node[circle, draw=prosvBlue, fill=white, minimum size=0.7cm] (a) at (-5,0) {A};
  \node[below, font=\scriptsize] at (-5,-0.5) {Алиса};
  \node[circle, draw=prosvBlue, fill=white, minimum size=0.7cm] (b) at (5,0) {B};
  \node[below, font=\scriptsize] at (5,-0.5) {Боб};
  
  \node[box] (enc) at (-2.3,0) {Шифрование $E$};
  \node[box] (dec) at (2.3,0) {Расшифрование $D$};
  
  \draw[->,>=stealth, prosvBlue, thick] (a) -- node[above]{$m$} (enc);
  \draw[->,>=stealth, prosvBlue, thick] (enc) -- node[above]{$c$} (dec);
  \draw[->,>=stealth, prosvBlue, thick] (dec) -- node[above]{$m$} (b);
  
  % Открытый ключ Боба
  \node[font=\Large] at (-2.3, 1.2) {[\textcolor{prosvGreen}{pk}]};
  \node[font=\scriptsize, text=prosvGreen, align=center] at (-2.3, 1.85) {откр. ключ\\ Боба $\mathit{pk}_B$};
  
  % Закрытый ключ Боба
  \node[font=\Large] at (2.3, 1.2) {[\textcolor{prosvRed}{sk}]};
  \node[font=\scriptsize, text=prosvRed, align=center] at (2.3, 1.85) {закр. ключ\\ Боба $\mathit{sk}_B$};
  
  % Стрелка от Боба к открытому ключу — "раздаёт всем"
  \draw[->,>=stealth, prosvGreen, dashed] (b) to[bend right=40] (-2.3, 1.4);
  \node[text=prosvGreen, font=\scriptsize] at (1.5, 2.5) {публикуется};
\end{tikzpicture}
\caption{Асимметричное шифрование: Алиса использует \emph{открытый} ключ Боба, который тот заранее опубликовал. Расшифровать может только Боб --- его \emph{закрытым} ключом.}
\label{fig:asymenc}
\end{figure}

Заметьте: больше нет проблемы распределения ключей. Алиса не должна заранее ничего секретного передавать Бобу --- она просто берёт его открытый ключ из открытого источника.

% -----------------------------------------------------------------
\subsection{Односторонние функции и функции с секретом}

В чём математическая суть асимметричной криптографии? В существовании \term{односторонних функций}.

\begin{definition}{}{def:p2-1}
\term{Односторонняя функция} (\engterm{one-way function}) --- это функция $f \colon X \to Y$, обладающая двумя свойствами:
\begin{itemize}
  \item \emph{Прямое вычисление просто:} зная $x$, легко (полиномиально быстро) вычислить $f(x)$.
  \item \emph{Обращение трудно:} зная $y = f(x)$, восстановить $x$ без какой-либо дополнительной информации --- задача \emph{практически} нерешаемая (требует экспоненциального времени).
\end{itemize}
\end{definition}

\begin{interestbox}
Существование «настоящих» односторонних функций математически не доказано (это связано с одной из главных нерешённых задач --- проблемой $\mathbf{P}\ne\mathbf{NP}$). На практике мы используем функции, для которых обращение \emph{считается} трудным, потому что человечество за десятилетия так и не научилось делать это быстро.
\end{interestbox}

Хороший бытовой образ: \emph{разбить вазу} --- легко, \emph{собрать обратно} --- очень тяжело. Или: \emph{перемешать колоду карт} --- мгновенная операция, \emph{восстановить исходный порядок} --- немыслимо.

Главные примеры из теории чисел:
\begin{itemize}
  \item \term{Умножение vs. факторизация:} перемножить $p\cdot q$ --- мгновенно; разложить произведение на простые --- трудно. Это база RSA.
  \item \term{Дискретный логарифм:} возвести $g$ в степень $x$ по модулю $p$ --- быстро; зная $g^x \bmod p$, восстановить $x$ --- очень трудно. Это база схемы Диффи--Хеллмана.
\end{itemize}

\paragraph*{Функции с секретом.} Просто односторонней функции для асимметричного шифрования мало: нужно, чтобы у \emph{кого-то одного} (владельца) была возможность всё-таки обратить функцию --- именно для расшифрования сообщений (рис.~\ref{fig:p2-4}).

\begin{definition}{}{def:p2-2}
\term{Функция с секретом} (\engterm{trapdoor function}) --- односторонняя функция $f$, для которой существует дополнительная секретная информация $t$ (\emph{лазейка}, \emph{trapdoor}), знание которой делает обращение $f$ полиномиально быстрым.
\end{definition}

\begin{figure}[H]
\centering
\begin{tikzpicture}[font=\small, every node/.style={align=center}]
  % Дверь без ключа
  \node[draw=prosvBlue, fill=prosvLight, rounded corners=2pt, minimum width=4cm, minimum height=1cm] (f) at (0,0) {Прямое вычисление:\\ \textit{очень быстро}};
  \node[draw=prosvRed, fill=prosvCream, rounded corners=2pt, minimum width=4cm, minimum height=1cm] (fi) at (0,-1.8) {Обращение без ключа:\\ \textit{экспоненциально долго}};
  \node[draw=prosvGreen, fill=prosvLight, rounded corners=2pt, minimum width=4cm, minimum height=1cm] (fis) at (0,-3.6) {Обращение со знанием секрета $t$:\\ \textit{полиномиально быстро}};
  
  \draw[->,>=stealth, thick, prosvBlue] (-3.0,0) -- (-2.05,0);
  \node[left, font=\scriptsize, text=prosvGray] at (-3,0) {$x$};
  \draw[->,>=stealth, thick, prosvBlue] (2.05,0) -- (3,0);
  \node[right, font=\scriptsize, text=prosvGray] at (3,0) {$f(x) = y$};
  
  \draw[->,>=stealth, thick, prosvRed] (3,-1.8) -- (2.05,-1.8);
  \node[right, font=\scriptsize, text=prosvGray] at (3,-1.8) {$y$};
  \draw[->,>=stealth, thick, prosvRed] (-2.05,-1.8) -- (-3,-1.8);
  \node[left, font=\scriptsize, text=prosvGray] at (-3,-1.8) {$x$};
  
  \draw[->,>=stealth, thick, prosvGreen] (3,-3.6) -- (2.05,-3.6);
  \node[right, font=\scriptsize, text=prosvGray] at (3,-3.6) {$y + t$};
  \draw[->,>=stealth, thick, prosvGreen] (-2.05,-3.6) -- (-3,-3.6);
  \node[left, font=\scriptsize, text=prosvGray] at (-3,-3.6) {$x$};
\end{tikzpicture}
\caption{Функция с секретом ведёт себя как односторонняя, но «открывается» при знании дополнительной информации~$t$.}
\label{fig:p2-4}
\end{figure}

Можно представить себе сейф с замком: запереть сейф (положить документ внутрь) может каждый, но открыть --- только владелец ключа.

\paragraph*{Идея использования.} Открытый ключ \emph{задаёт} функцию с секретом $f$. Кто угодно может вычислить $c = f(m)$ --- зашифровать сообщение. Для расшифрования нужно обратить $f$, что без секрета $t$ --- безнадёжно. У владельца секрета $t$ это вычисление быстрое. Тем самым роль закрытого ключа --- хранить лазейку $t$.

% -----------------------------------------------------------------
\subsection{Где криптография работает прямо сейчас}

Чтобы дальнейший материал не казался академическим упражнением, перечислим, в каких знакомых каждому ситуациях работают (часто прямо сейчас, пока вы это читаете) описанные в этой главе идеи.

\begin{itemize}
  \item \textbf{Загрузка любого сайта} (\texttt{https}). Браузер договаривается с сервером о ключе через схему Диффи--Хеллмана, после чего весь трафик шифруется симметричным алгоритмом (\engterm{AES}).
  \item \textbf{Мессенджеры} (WhatsApp, Signal, Telegram). WhatsApp и Signal по умолчанию используют \emph{сквозное шифрование} (\engterm{end-to-end}): сервер видит лишь шифр-текст. В Telegram сквозное шифрование включается только в «секретных чатах», а обычные облачные чаты шифруются лишь между клиентом и сервером (и хранятся на серверах Telegram). Согласование ключей --- асимметричное.
  \item \textbf{Банковские карты и оплата.} Транзакция от карты проходит через цепочку обмена цифровыми подписями.
  \item \textbf{Госуслуги и налоговая.} Подача декларации сопровождается \term{электронной цифровой подписью} (ЭЦП), о которой подробно --- в §~\ref{sec:rsa-apps}.
  \item \textbf{Цифровой пропуск/QR-код} (от билета в кино до подтверждения вакцинации). Содержит подпись эмитента, проверяемую его открытым ключом.
  \item \textbf{Криптовалюты.} Биткойн и его «родственники» --- огромная распределённая система с электронными подписями: кошелёк --- это просто пара (открытый ключ, закрытый ключ).
\end{itemize}

\paragraph*{Что должно быть у «правильной» криптосистемы.} Когда мы строим конкретную схему, надо проверять, что она удовлетворяет следующим естественным требованиям:
\begin{enumerate}
  \item \emph{Корректность.} Если Алиса честно зашифровала открытым ключом Боба, Боб всегда расшифрует --- ровно исходное сообщение.
  \item \emph{Стойкость.} Подслушивающий, видя только шифр-текст и открытый ключ, не может практически восстановить исходное сообщение.
  \item \emph{Эффективность.} Шифрование и расшифрование выполняются быстро (полиномиально).
  \item \emph{Защита ключа.} Закрытый ключ не должен «утекать» из открытого --- даже зная открытый ключ, нельзя восстановить закрытый.
\end{enumerate}

В следующих параграфах мы построим конкретные схемы --- RSA и Диффи--Хеллман --- и убедимся, что они этим требованиям удовлетворяют.

\begin{tasksbox}
\begin{enumerate}
  \item Зашифруйте шифром Цезаря со сдвигом 5 слово \texttt{INFORMATIKA}.
  \item Почему симметричное шифрование плохо подходит для миллиарда пользователей? Подсчитайте, сколько ключей надо хранить.
  \item Объясните своими словами, в чём состоит свойство «односторонности» функции.
  \item* Знаменитая Энигма имела примерно $10^{20}$ возможных стартовых установок. Хватило бы перебора, если перебирать миллион вариантов в секунду?
  \item Приведите три бытовые операции, которые ведут себя как «односторонние функции».
\end{enumerate}
\end{tasksbox}

% =============================================================
%  § 3.  Криптосистемы RSA и Диффи--Хеллмана
% =============================================================

\section{Криптосистемы RSA и Диффи--Хеллмана}
\label{sec:rsa-dh}

В этом параграфе мы построим две настоящие, работающие сегодня схемы асимметричной криптографии. Обе появились почти в одно время --- между 1976 и 1978 годами --- и обе опираются на результаты предыдущего параграфа.

\subsection{Криптосистема RSA}

Идея была опубликована в 1978 году тремя сотрудниками Массачусетского технологического института --- \emph{Р}ональдом Райвестом, \emph{Ш}амиром (Ади Шамир) и \emph{А}длеманом (Леонард Адлеман). Их фамилии и дали название системе --- RSA.

\paragraph*{Генерация ключей.} Каждый пользователь, желающий получать зашифрованные сообщения, выполняет одноразовую процедуру.

\begin{algorithmbox}[title={Алгоритм: генерация ключей RSA}]
\begin{enumerate}
  \item Выбрать два больших случайных \emph{различных} простых числа $p$ и $q$ (на практике --- по 1024 бита каждое).
  \item Вычислить $n = p\cdot q$ и $\Eul(n) = (p-1)(q-1)$.
  \item Выбрать целое $e$, взаимно простое с $\Eul(n)$ (часто берут $e = 65537$).
  \item С помощью расширенного алгоритма Евклида найти $d$ такое, что
  \[
  e\cdot d \equiv 1 \modn{\Eul(n)}.
  \]
  \item \term{Открытый ключ:} пара $(n, e)$ --- публикуется.\\
        \term{Закрытый ключ:} число $d$ --- хранится в секрете.\\
        \term{Лишние данные:} числа $p, q, \Eul(n)$ безопасно уничтожаются.
\end{enumerate}
\end{algorithmbox}

\paragraph*{Шифрование.} Алиса хочет послать сообщение $m$ Бобу. Сообщение представлено числом из отрезка $[0, n-1]$ (длинное сообщение разбивается на блоки). Алиса берёт открытый ключ Боба $(n, e)$ и вычисляет:
\[
c \;=\; m^e \bmod n.
\]
Число $c$ называется \term{шифр-текстом} и передаётся открыто.

\paragraph*{Расшифрование.} Боб, получив $c$, использует свой закрытый ключ $d$ и вычисляет:
\[
m' \;=\; c^d \bmod n.
\]
Утверждается: $m' = m$ (рис.~\ref{fig:rsa-scheme}).

\begin{figure}[H]
\centering
\begin{tikzpicture}[
  font=\small,
  every node/.style={align=center},
  box/.style={rectangle, draw=prosvBlue, fill=prosvLight,
              minimum width=3.0cm, minimum height=1.1cm,
              rounded corners=2pt, font=\small}
]
  \node[circle, draw=prosvBlue, fill=white, minimum size=0.8cm] (a) at (-5.5,0) {A};
  \node[below, font=\scriptsize] at (-5.5,-0.55) {Алиса};
  \node[circle, draw=prosvBlue, fill=white, minimum size=0.8cm] (b) at (5.5,0) {B};
  \node[below, font=\scriptsize] at (5.5,-0.55) {Боб};
  
  \node[box] (enc) at (-2.5,0) {$c = m^e \bmod n$};
  \node[box] (dec) at (2.5,0) {$m = c^d \bmod n$};
  
  \draw[->,>=stealth, prosvBlue, thick] (a) -- node[above]{$m$} (enc);
  \draw[->,>=stealth, prosvBlue, thick] (enc) -- node[above]{$c$} (dec);
  \draw[->,>=stealth, prosvBlue, thick] (dec) -- node[above]{$m$} (b);
  
  % Ключи
  \node[font=\scriptsize, text=prosvGreen, align=center] at (-2.5, 1.4) {откр. ключ Боба\\$(n, e)$};
  \node[font=\scriptsize, text=prosvRed, align=center] at (2.5, 1.4) {закр. ключ Боба\\$d$};
  
  \draw[->, dashed, prosvGreen] (-2.5, 1.0) -- (-2.5, 0.6);
  \draw[->, dashed, prosvRed] (2.5, 1.0) -- (2.5, 0.6);
  
  \draw[->,>=stealth, prosvGreen, dashed] (b) to[bend right=50] (-2.5, 1.6);
  \node[font=\scriptsize, text=prosvGreen] at (1, 2.4) {опубликован};
\end{tikzpicture}
\caption{Схема работы RSA: шифрование --- возведение в степень $e$ по модулю $n$, расшифрование --- возведение в степень $d$. Обе операции эффективны благодаря быстрому возведению в степень.}
\label{fig:rsa-scheme}
\end{figure}

% -----------------------------------------------------------------
\subsection{Почему RSA работает: корректность}

\begin{theorem}{корректность RSA}{thm:p3-1}
В обозначениях процедуры RSA, для любого $m \in [0, n-1]$ выполняется
\[
(m^e)^d \equiv m \modn{n}.
\]
\end{theorem}

\begin{proof}
По построению $e\cdot d \equiv 1 \modn{\Eul(n)}$, значит, найдётся целое $k$, такое что $e\cdot d = 1 + k\cdot \Eul(n)$. Подставим:
\[
m^{ed} = m^{1 + k\Eul(n)} = m\cdot \bigl(m^{\Eul(n)}\bigr)^k.
\]

\textbf{Случай 1.} $\NOD(m, n) = 1$. Тогда по теореме Эйлера $m^{\Eul(n)} \equiv 1 \modn n$, и
\[
m^{ed} \equiv m\cdot 1^k = m \modn n.
\]

\textbf{Случай 2.} $\NOD(m, n) > 1$. Поскольку $n = pq$ и $p, q$ простые, это значит, что $m$ делится либо на $p$, либо на $q$ (но не на оба --- иначе $m \ge pq = n$, чего быть не может). Рассмотрим, например, $p \mid m$. Тогда $m \equiv 0 \modn p$, и обе части сравнения $m^{ed} \equiv m \modn p$ обращаются в ноль --- утверждение выполнено.

С другой стороны, $\NOD(m, q) = 1$, поэтому по малой теореме Ферма $m^{q-1} \equiv 1 \modn q$. А значит, $m^{(p-1)(q-1)} = m^{\Eul(n)} \equiv 1 \modn q$, и
\[
m^{ed} = m\cdot \bigl(m^{\Eul(n)}\bigr)^k \equiv m\cdot 1^k = m \modn q.
\]
Получили $m^{ed} \equiv m$ одновременно по модулям $p$ и $q$. По китайской теореме об остатках (или просто потому, что разность $m^{ed} - m$ делится и на $p$, и на $q$, а значит, на их произведение $n$) сравнение $m^{ed} \equiv m \modn n$ выполнено. \qedhere
\end{proof}

\paragraph*{Учебный пример (маленькие числа).} Полноценный RSA использует тысячебитные числа, но для понимания работает и крошечный пример.

\begin{example}{генерация и применение ключа RSA}{ex:p3-1}
Пусть $p = 11$, $q = 23$. Тогда:
\begin{itemize}
  \item $n = 11\cdot 23 = 253$, $\Eul(n) = 10\cdot 22 = 220$.
  \item Выберем $e = 7$. Проверим взаимную простоту: $\NOD(7, 220) = 1$.
  \item Найдём $d = 7^{-1} \bmod 220$. Расширенным алгоритмом Евклида:
  \[
  220 = 31\cdot 7 + 3,\quad 7 = 2\cdot 3 + 1.
  \]
  Раскручиваем: $1 = 7 - 2\cdot 3 = 7 - 2(220 - 31\cdot 7) = 63\cdot 7 - 2\cdot 220$. Значит, $d = 63$ (проверка: $7\cdot 63 = 441 = 2\cdot 220 + 1$).
  \item Открытый ключ: $(253, 7)$. Закрытый ключ: $d = 63$.
\end{itemize}

\emph{Шифрование.} Пусть $m = 50$. Считаем $c = 50^7 \bmod 253$ быстрым возведением:
\[
50^2 = 2500 \equiv 2500 - 9\cdot 253 = 223,\qquad 50^4 \equiv 223^2 = 49\,729.
\]
Делим $49\,729 / 253 \approx 196{,}5$, $196\cdot 253 = 49\,588$, остаток $141$. Значит, $50^4 \equiv 141$. Далее $50^7 = 50^4\cdot 50^2\cdot 50 \equiv 141\cdot 223\cdot 50 \bmod 253$. Считаем: $141\cdot 223 = 31\,443 = 124\cdot 253 + 71$, значит $\equiv 71$. Затем $71\cdot 50 = 3550 = 14\cdot 253 + 8$, значит $\equiv 8$. Итак, $c = 8$.

\emph{Расшифрование.} Боб считает $c^d = 8^{63} \bmod 253$. Вместо длинных вычислений «вручную» можно использовать малую теорему Ферма по модулям 11 и 23 раздельно, а затем китайскую теорему об остатках --- именно так и поступают на практике. Ответ: $50$.
\end{example}

\begin{interestbox}
Удивительный исторический парадокс: ровно та же схема была независимо открыта в 1973 году британским математиком Клиффордом Коксом, работавшим в Британском управлении правительственной связи (GCHQ). Но работа была засекречена и рассекречена только в 1997 году. К этому моменту Райвест, Шамир и Адлеман уже стали всемирно известными, основали компанию RSA Security, а Шамир получил Премию Тьюринга. Типичный пример того, как секретность тормозит развитие науки.
\end{interestbox}

% -----------------------------------------------------------------
\subsection{Почему RSA безопасен}

\begin{importantbox}
Стойкость RSA опирается на предположение: \emph{задача факторизации $n$ практически нерешаема при больших $n$}.
\end{importantbox}

Действительно, если злоумышленник умеет факторизовать $n$, он восстанавливает $p, q$, затем $\Eul(n) = (p-1)(q-1)$ и наконец $d = e^{-1} \bmod \Eul(n)$ --- получит закрытый ключ. Обратное (что взлом RSA эквивалентен факторизации) также считается верным, но строго не доказано: это одна из открытых проблем в криптографии.

Заметим, что \emph{знать $n$ и $e$} (открытый ключ) --- ещё далеко не значит знать $\Eul(n)$. Знание $\Eul(n)$ и знание разложения $n = pq$ --- одна и та же информация: ведь $p + q = n - \Eul(n) + 1$, а $p - q = \sqrt{(p+q)^2 - 4n}$. То есть «нахождение $\Eul(n)$» и «факторизация $n$» --- задачи одинаковой сложности.

\paragraph*{Почему нужны именно \emph{большие} простые числа?} Современный рекомендуемый размер $n$ --- 2048 или даже 4096 бит. Чтобы взломать 2048-битный RSA, нужно факторизовать число длиной около 617 десятичных цифр; лучший известный алгоритм потратит на это больше времени, чем существует Вселенная.

\begin{example}{границы практической факторизации}{ex:p3-2}
В 2020 году группе исследователей удалось факторизовать число RSA-250 (250 десятичных цифр $\approx$ 829 бит). Заняло это около 2700 ядро-лет вычислений --- то есть один компьютер с 2700 ядрами работал бы целый год. Это подтверждает, что 1024-битный RSA скоро будет «на грани», а 2048-битный --- надолго безопасен.
\end{example}

\paragraph*{Проверим выполнение требований из § \ref{sec:history}.}
\begin{enumerate}
  \item \emph{Корректность} --- доказана теоремой выше.
  \item \emph{Стойкость} --- опирается на сложность факторизации.
  \item \emph{Эффективность} --- шифрование/расшифрование сводятся к возведению в степень по модулю, что мы умеем делать полиномиально (см. §~\ref{subsec:fastpow}).
  \item \emph{Защита ключа} --- знание $e, n$ не даёт быстрого способа найти $d$ без факторизации.
\end{enumerate}

% -----------------------------------------------------------------
\subsection{Протокол Диффи--Хеллмана}

Уитфилд Диффи и Мартин Хеллман решали несколько другую (хотя и родственную) задачу: \emph{договориться} о секретном ключе по открытому каналу, --- не передавая по нему ничего секретного. Их работа 1976 года так и называется: «Новые направления в криптографии».

\paragraph*{Идея на пальцах: краски.} Прежде чем вводить формулы, разберём наглядную аналогию (рис.~\ref{fig:p3-2}).

\begin{figure}[H]
\centering
\begin{tikzpicture}[
  font=\small, every node/.style={align=center},
  circ/.style={circle, draw, minimum size=0.7cm}
]
  % Алиса и Боб + общая краска
  \node[circ, fill=yellow!50!brown!20] (g) at (0,3) {};
  \node[below, font=\scriptsize] at (0,2.6) {общая «жёлтая»};
  
  \node[circ, fill=red!60!brown!50] (sa) at (-5,3) {};
  \node[below, font=\scriptsize, align=center] at (-5,2.6) {Алисин «красный»\\(секрет)};
  
  \node[circ, fill=blue!60!cyan!50] (sb) at (5,3) {};
  \node[below, font=\scriptsize, align=center] at (5,2.6) {Бобов «синий»\\(секрет)};
  
  % Смеси
  \node[circ, fill=orange!70!red!30] (ya) at (-2.5,1.0) {};
  \node[below, font=\scriptsize, align=center] at (-2.5,0.6) {жёлт. $+$ красн.\\(передаётся\\открыто)};
  
  \node[circ, fill=green!60!blue!30] (yb) at (2.5,1.0) {};
  \node[below, font=\scriptsize, align=center] at (2.5,0.6) {жёлт. $+$ син.\\(передаётся\\открыто)};
  
  \draw[->, prosvBlue] (g) -- (ya);
  \draw[->, prosvBlue] (sa) -- (ya);
  \draw[->, prosvBlue] (g) -- (yb);
  \draw[->, prosvBlue] (sb) -- (yb);
  
  % Итоговая общая смесь
  \node[circ, fill=brown!50] (final) at (0,-2.0) {};
  \node[below, font=\scriptsize, align=center] at (0,-2.4) {итоговая «коричневая»\\ --- общий секрет};
  
  \draw[->, prosvGreen, thick] (yb) -- node[above right, font=\scriptsize, text=prosvGreen]{+ Алисин «красный»} (final);
  \draw[->, prosvGreen, thick] (ya) -- node[above left, font=\scriptsize, text=prosvGreen]{+ Бобов «синий»} (final);
\end{tikzpicture}
\caption{Аналогия Диффи--Хеллмана со смешиванием красок. Смешать --- легко, разделить смесь обратно на исходные цвета --- невозможно. Алиса и Боб приходят к одному цвету разными путями.}
\label{fig:p3-2}
\end{figure}

\noindent
Алиса и Боб условились о «жёлтой» базовой краске (известна всем). Каждый выбрал свою секретную краску (Алиса --- красный, Боб --- синий) и публично передал смесь «своя $+$ жёлтая». Подслушивающий видит обе смеси, но не может разложить их обратно. А Алиса и Боб, добавив к чужой смеси свою секретную краску, приходят к одинаковому итогу: жёлт. $+$ красн. $+$ син.

\paragraph*{Математическая реализация.} Роль «смешивания» играет \term{дискретное возведение в степень} по модулю простого числа (рис.~\ref{fig:dh-scheme}).

\begin{algorithmbox}[title={Алгоритм: протокол Диффи--Хеллмана}]
\textbf{Подготовка (общедоступно):} большое простое число $p$ и целое $g$, $1 < g < p$.
\begin{enumerate}
  \item Алиса выбирает случайный секрет $a \in [1, p-2]$ и посылает Бобу $A = g^a \bmod p$.
  \item Боб выбирает случайный секрет $b \in [1, p-2]$ и посылает Алисе $B = g^b \bmod p$.
  \item Алиса вычисляет общий ключ: $K = B^a \bmod p$.
  \item Боб вычисляет общий ключ: $K = A^b \bmod p$.
\end{enumerate}
\textbf{Результат.} Алиса и Боб получили один и тот же ключ
\[
K = g^{ab} \bmod p,
\]
который никогда не передавался по каналу связи.
\end{algorithmbox}

\begin{figure}[H]
\centering
\begin{tikzpicture}[
  font=\small,
  every node/.style={align=center},
  alice/.style={fill=prosvLight, draw=prosvBlue, rounded corners=2pt, minimum width=4.5cm, minimum height=0.7cm},
  bob/.style={fill=prosvCream, draw=prosvGold!80!brown, rounded corners=2pt, minimum width=4.5cm, minimum height=0.7cm}
]
  % Заголовки
  \node[font=\bfseries\color{prosvBlue}] at (-3.5, 4) {Алиса};
  \node[font=\bfseries\color{prosvRed}] at (3.5, 4) {Боб};
  \node[font=\bfseries\color{prosvGreen}] at (0, 4) {Открытый канал};
  
  % Шаги Алисы
  \node[alice] (A1) at (-3.5, 3) {Выбирает секрет $a$};
  \node[alice] (A2) at (-3.5, 2) {Считает $A = g^a \bmod p$};
  \node[alice] (A3) at (-3.5, 0.5) {Получает $B$};
  \node[alice] (A4) at (-3.5, -0.5) {$K = B^a \bmod p$};
  
  % Шаги Боба
  \node[bob] (B1) at (3.5, 3) {Выбирает секрет $b$};
  \node[bob] (B2) at (3.5, 2) {Считает $B = g^b \bmod p$};
  \node[bob] (B3) at (3.5, 0.5) {Получает $A$};
  \node[bob] (B4) at (3.5, -0.5) {$K = A^b \bmod p$};
  
  % Передача
  \draw[->, thick, prosvBlue] (A2) -- node[above, font=\scriptsize]{$A$} (B3);
  \draw[->, thick, prosvRed]  (B2) -- node[below, font=\scriptsize]{$B$} (A3);
  
  % Итог
  \node[fill=prosvLight!30!white, draw=prosvGreen, very thick, rounded corners=3pt, minimum width=7cm, minimum height=0.8cm] at (0,-1.8) {Общий ключ $K = g^{ab} \bmod p$};
\end{tikzpicture}
\caption{Диаграмма обмена в протоколе Диффи--Хеллмана.}
\label{fig:dh-scheme}
\end{figure}

\paragraph*{Учебный пример.}

\begin{example}{обмен ключами Диффи--Хеллмана}{ex:p3-3}
Пусть $p = 23$, $g = 5$.
\begin{itemize}
  \item Алиса выбирает $a = 6$, отправляет $A = 5^6 \bmod 23$. Считая последовательно по таблице степеней пятёрки:
  \[
  5^1\equiv 5,\;5^2\equiv 2,\;5^3\equiv 10,\;5^4\equiv 4,\;5^5\equiv 20,\;5^6\equiv 8,
  \]
  получаем $A = 8$.
  \item Боб выбирает $b = 15$, отправляет $B = 5^{15} \bmod 23$. Продолжая таблицу:
  \[
  5^7\equiv 17,\;5^8\equiv 16,\;5^9\equiv 11,\;5^{10}\equiv 9,\;5^{11}\equiv 22,
  \]
  \[
  5^{12}\equiv 18,\;5^{13}\equiv 21,\;5^{14}\equiv 13,\;5^{15}\equiv 19,
  \]
  получаем $B = 19$.
  \item Алиса вычисляет общий ключ $K = B^a = 19^6 \bmod 23$. Замечая, что $19 \equiv -4 \modn{23}$, имеем $(-4)^6 = 4^6 = 4096 = 178\cdot 23 + 2$, значит, $K = 2$.
  \item Боб вычисляет $K = A^b = 8^{15} \bmod 23$. Но $8 = 5^6$, значит, $8^{15} = 5^{90}$. По малой теореме Ферма $5^{22} \equiv 1 \modn{23}$, а $90 = 4\cdot 22 + 2$, поэтому $5^{90} \equiv 5^2 = 2 \modn{23}$. Получаем $K = 2$.
\end{itemize}

Согласование удалось: общий ключ $K = 2$. И это притом, что числа $a = 6$ и $b = 15$ никогда не передавались по каналу!
\end{example}

\paragraph*{Безопасность Диффи--Хеллмана.} Чтобы подслушивающий, видя $g, p, A, B$, мог найти общий ключ $K = g^{ab}$, ему нужно сначала восстановить хотя бы один из секретов $a$ или $b$, --- то есть решить уравнение
\[
g^a \equiv A \modn p
\]
относительно $a$ при известных $g, A, p$. Эта задача называется \term{задачей дискретного логарифма} (в группе $\Zn{p}^*$).

\begin{importantbox}
Для всех известных классических алгоритмов задача дискретного логарифма по модулю большого простого --- субэкспоненциальная (примерно той же сложности, что и факторизация). Однако \emph{по модулю эллиптической кривой} --- значительно сложнее, и поэтому современные мессенджеры обычно используют вариант Диффи--Хеллмана не по модулю числа, а на эллиптических кривых (\engterm{ECDH}). Идея та же.
\end{importantbox}

% -----------------------------------------------------------------
\subsection{Атака «человек посередине»}

И RSA, и Диффи--Хеллман в чистом виде уязвимы к одной атаке. Представим, что между Алисой и Бобом «вклинился» злоумышленник --- Мэллори (\engterm{Man-in-the-Middle}) (рис.~\ref{fig:p3-4}).

\begin{figure}[H]
\centering
\begin{tikzpicture}[font=\small, every node/.style={align=center}]
  \node[circle, draw=prosvBlue, fill=white, minimum size=0.8cm] (a) at (-5,0) {A};
  \node[below, font=\scriptsize] at (-5,-0.55) {Алиса};
  
  \node[circle, draw=prosvRed, fill=prosvCream, minimum size=0.8cm] (m) at (0,0) {M};
  \node[below, font=\scriptsize, text=prosvRed] at (0,-0.55) {Мэллори};
  
  \node[circle, draw=prosvBlue, fill=white, minimum size=0.8cm] (b) at (5,0) {B};
  \node[below, font=\scriptsize] at (5,-0.55) {Боб};
  
  \draw[->, prosvBlue, thick] (a) -- node[above, font=\scriptsize]{«Боб, мой $A$»} (m);
  \draw[->, prosvRed, thick]  (m) -- node[above, font=\scriptsize]{«Я Алиса, мой $A'$»} (b);
  \draw[<-, prosvBlue, thick] (a) to[bend right=30] node[below, font=\scriptsize]{«моё $B''$»} (m);
  \draw[<-, prosvRed, thick]  (m) to[bend right=30] node[below, font=\scriptsize]{«моё $B$»} (b);
\end{tikzpicture}
\caption{Атака «человек посередине»: Мэллори создаёт две независимые сессии, в каждой выдавая себя за противоположную сторону.}
\label{fig:p3-4}
\end{figure}

Мэллори перехватывает $A$ от Алисы, заменяет на свой $A'$ и пересылает Бобу. С Бобом он устанавливает один ключ, с Алисой --- другой. Затем расшифровывает каждое сообщение, читает (и при желании меняет), снова шифрует и передаёт дальше. Внешне Алиса и Боб ничего не замечают.

\paragraph*{Защита.} Нужна \emph{аутентификация} --- способ убедиться, что собеседник действительно тот, за кого себя выдаёт. В реальных системах это решается через \term{цифровые сертификаты} и \term{электронную цифровую подпись}. О них --- следующий параграф.

% -----------------------------------------------------------------
\begin{tasksbox}
\begin{enumerate}
  \item В системе RSA дано $p = 7$, $q = 13$, $e = 5$. Найдите $n$, $\Eul(n)$ и $d$.
  \item Используя ключ из задачи 1, зашифруйте сообщение $m = 9$.
  \item В Диффи--Хеллмане положите $p = 17$, $g = 3$, $a = 5$, $b = 7$. Найдите $A$, $B$ и общий ключ $K$.
  \item Почему в RSA нельзя брать $e = 1$? А $e = \Eul(n) - 1$?
  \item Что произойдёт, если в RSA Алиса случайно использует $m \ge n$?
  \item* Объясните, как именно атака «человек посередине» работает в случае RSA (не Диффи--Хеллмана).
  \item* Что вы посоветуете человеку, который хочет «шифровать всё своими силами без сертификатов»? Какие угрозы он не учитывает?
\end{enumerate}
\end{tasksbox}

% =====================================================================
%   §4. Применения RSA: аутентификация, цифровая подпись, голосование
% =====================================================================
\section{Применения RSA: аутентификация, цифровая подпись, электронное голосование}
\label{sec:rsa-apps}

В предыдущем параграфе мы научились с помощью RSA \emph{передавать секретное сообщение}. На первый взгляд кажется, что этого уже достаточно для всех задач сетевой безопасности. Однако в действительности шифрование решает лишь одну из проблем --- \term{конфиденциальность}. Помимо неё, в современных системах необходимо обеспечивать ещё как минимум три:

\begin{itemize}
  \item \term{Аутентификацию} --- уверенность, что вы общаетесь именно с тем, за кого себя выдаёт собеседник.
  \item \term{Целостность} --- уверенность, что сообщение не было незаметно изменено по дороге.
  \item \term{Неотказуемость} --- невозможность для отправителя впоследствии отказаться от своего сообщения («это не я отправлял»).
\end{itemize}

Удивительно, но все эти задачи можно решить с помощью \emph{той же самой схемы RSA}, если применить её не «прямо», а наоборот --- сначала закрытым ключом, а потом открытым. Об этом и пойдёт речь в этом параграфе.

% ---------------------------------------------------------------------
\subsection{Аутентификация: «докажи, что это ты»}
\label{subsec:auth}

Представим типичную ситуацию: Алиса заходит в свой личный кабинет на сайте банка. Банк хочет убедиться, что это действительно его клиентка, а не злоумышленник, подсмотревший пароль.

Самый простой подход --- запросить пароль. Но у него есть очевидные недостатки: пароль можно подсмотреть, перехватить, угадать, заставить пользователя его выдать. Кроме того, чтобы проверить пароль, банк должен \emph{хранить} его (или его хеш --- см. §\ref{sec:hashing}) у себя на серверах. А значит, при взломе сервера утечёт сразу множество паролей.

Криптография позволяет построить намного более изящную схему, в которой \emph{секрет вообще никогда не передаётся по сети}.

\begin{importantbox}[title={Важно: принцип «доказательства с нулевым разглашением»}]
Сторона $A$ должна доказать стороне $B$, что обладает неким секретом, \emph{не выдавая никакой информации о самом секрете}. Идея, кажущаяся парадоксальной, на практике вполне реализуема и лежит в основе многих криптографических протоколов.
\end{importantbox}

Простейший пример --- \term{протокол challenge-response} («запрос--ответ») на базе RSA. Пусть у Алисы заведена пара ключей: открытый $(e, n)$ опубликован в системе банка, а закрытый $d$ хранится только у неё (например, в специальном USB-токене или в защищённом хранилище смартфона).

\begin{algorithmbox}[title={Алгоритм: аутентификация через RSA}]
\begin{enumerate}
  \item \textbf{Запрос (challenge).} Банк генерирует \emph{случайное} число $r$ ($0 < r < n$), запоминает его и отсылает Алисе.
  \item \textbf{Ответ (response).} Алиса вычисляет $s = r^{d} \bmod n$, используя свой закрытый ключ, и отправляет $s$ обратно.
  \item \textbf{Проверка.} Банк вычисляет $r' = s^{e} \bmod n$ и сравнивает с тем $r$, которое он отсылал. Если $r' = r$, проверка пройдена.
\end{enumerate}
\end{algorithmbox}

\paragraph*{Почему это работает?} В предыдущем параграфе мы доказали, что для любого $r$ выполняется
\[
  (r^{d})^{e} \equiv r \pmod n.
\]
То есть только тот, кто знает $d$, может «правильно» преобразовать $r$ так, чтобы после открытого преобразования $\cdot^{e}$ получилось снова $r$. Любой другой пользователь, не знающий $d$, может лишь \emph{угадывать} ответ --- а вероятность угадать одно нужное число из миллиарда миллиардов исчезающе мала.

\paragraph*{Почему это \emph{безопасно}?}
\begin{itemize}
  \item Алиса \emph{не передаёт} свой закрытый ключ $d$ --- он остаётся у неё.
  \item Каждый раз $r$ \emph{случайное и новое}. Поэтому даже если злоумышленник перехватит весь сеанс и услышит пару $(r, s)$, в следующий раз банк пришлёт другое $r$, и старый ответ $s$ не подойдёт. Атака повторного воспроизведения (replay attack) исключена.
  \item Банк \emph{не хранит секретов клиента}: у него лишь открытый ключ, утечка которого ничем не страшна.
\end{itemize}

\begin{importantbox}[title={Важно: это учебная иллюстрация}]
Описанная схема передаёт \emph{идею} challenge-response, но не является полноценным протоколом с нулевым разглашением. В реальных системах нельзя позволять кому угодно получать $r^{d} \bmod n$ для произвольного $r$: устройство фактически превращается в «оракул» RSA-подписи. Поэтому на практике запрос специально кодируют, а ключи для подписи и аутентификации разделяют.
\end{importantbox}

Схему этого протокола в графическом виде даёт рис.~\ref{fig:p4-1}.

\begin{figure}[H]
\centering
\begin{tikzpicture}[
  >=Stealth,
  every node/.style={font=\small},
  party/.style={draw=prosvBlue, rounded corners, fill=prosvLight, minimum width=2.2cm, minimum height=1.1cm, align=center, thick}
]
  \node[party] (a) at (0,0) {\textbf{Алиса}\\ ключ $d$};
  \node[party] (b) at (8,0) {\textbf{Банк}\\ ключ $e$};

  \draw[->, thick, prosvRed] ([yshift=15pt]b.west) -- node[above, font=\footnotesize]{случайный $r$} ([yshift=15pt]a.east);
  \draw[<-, thick, prosvGreen!70!black] ([yshift=-5pt]b.west) -- node[below, font=\footnotesize]{$s = r^d \bmod n$} ([yshift=-5pt]a.east);

  \node[align=center, font=\footnotesize] at (4, -1.4) {Банк проверяет: $s^e \bmod n \stackrel{?}{=} r$};
\end{tikzpicture}
\caption{Схема аутентификации по запросу--ответу.}
\label{fig:p4-1}
\end{figure}

\begin{interestbox}[title={Это интересно: «запомни всё, чтобы ничего не помнить»}]
В современных смартфонах закрытый ключ $d$ хранится в специальном чипе --- \emph{защищённом анклаве} (Secure Enclave у Apple, Trusted Execution Environment на Android). Этот чип \emph{физически не позволяет считать} закрытый ключ извне. На него можно лишь «передать» число и попросить выполнить с ним операцию $\cdot^d \bmod n$. Так устроены, например, биометрические разблокировки --- отпечаток пальца «разрешает» чипу применить ключ, но не извлекает его.
\end{interestbox}

% ---------------------------------------------------------------------
\subsection{Электронная цифровая подпись}
\label{subsec:eds}

Аутентификация решает задачу «докажи, что ты --- это ты». Теперь рассмотрим родственную, но другую: «докажи, что \emph{именно ты} отправил \emph{именно это} сообщение». Эта задача стоит особенно остро при подписании документов: договоров, заявлений, налоговых деклараций.

Бумажная подпись от руки выполняет сразу две функции: \emph{идентифицирует} автора и \emph{связывает} его с конкретным документом. К сожалению, при простом копировании в электронном виде это свойство теряется: цифровое изображение подписи легко вырезать и приклеить к любому документу. Нужна принципиально иная конструкция.

\begin{definition}{электронная цифровая подпись}{def:p4-1}
\term{Электронная цифровая подпись} (ЭЦП) --- это короткий блок данных $\sigma$, привязываемый к документу $m$, обладающий двумя свойствами:
\begin{itemize}
  \item только владелец секретного ключа может сгенерировать $\sigma$ для произвольного $m$;
  \item любой желающий, зная открытый ключ, может проверить, что $\sigma$ действительно соответствует $m$ и подписан данным владельцем.
\end{itemize}
\end{definition}

Идея конструкции на основе RSA удивительно проста и красива: достаточно поменять местами роли ключей в шифровании.

\begin{algorithmbox}[title={Алгоритм: цифровая подпись RSA}]
Алиса имеет открытый ключ $(e, n)$ и закрытый $d$. Хочет подписать документ $m$ ($0 < m < n$).
\begin{enumerate}
  \item \textbf{Подпись.} Алиса вычисляет
  \[
    \sigma = m^{d} \bmod n
  \]
  и публикует пару $(m, \sigma)$.
  \item \textbf{Проверка.} Любой проверяющий вычисляет $m' = \sigma^{e} \bmod n$ и сравнивает с $m$. Если $m' = m$, подпись подлинна.
\end{enumerate}
\end{algorithmbox}

\paragraph*{Почему это подпись?} По тому же самому соотношению $(m^{d})^{e} \equiv m \pmod n$. Только владелец закрытого ключа $d$ может построить такое $\sigma$, что после возведения в степень $e$ получится исходное $m$. Если бы это умел кто-то другой, он, по сути, умел бы разлагать $n$ на множители --- а это, как мы обсуждали, считается практически невыполнимой задачей.

\paragraph*{Привязка к документу.} В отличие от рукописной подписи, ЭЦП \emph{зависит от содержимого}: если изменить хоть один бит в $m$, изменится и $\sigma$. Невозможно «отделить» подпись от документа и приклеить её к другому: для другого документа потребуется заново вычислять $\cdot^{d} \bmod n$, что без $d$ невыполнимо (рис.~\ref{fig:p4-2}).

\begin{figure}[H]
\centering
\begin{tikzpicture}[
  >=Stealth, every node/.style={font=\small},
  doc/.style={draw=prosvBlue, fill=prosvCream, rounded corners, minimum width=2cm, minimum height=1cm, align=center},
  op/.style={draw=prosvRed, fill=prosvLight, rounded corners, minimum width=2.6cm, minimum height=0.8cm, align=center, thick},
]
  \node[doc] (m1) at (0,0) {Документ $m$};
  \node[op]  (sign) at (3.7,0) {$\sigma = m^d \bmod n$};
  \node[doc] (pair) at (8,0) {$(m,\sigma)$};
  \draw[->, thick] (m1) -- (sign);
  \draw[->, thick] (sign) -- (pair);

  \node[op]  (verify) at (3.7,-2) {$\sigma^e \bmod n$};
  \node[doc] (mres) at (8,-2)   {$m'$};
  \node[align=center, font=\footnotesize] at (8.8,-3) {$m' \stackrel{?}{=} m$};

  \draw[->, thick] (pair.south) -- ++(0,-0.6) -| (verify.north);
  \draw[->, thick] (verify) -- (mres);

  \node[align=left, font=\footnotesize\itshape, prosvBlue] at (-0.5,-2) {Подпись (закр.\ ключ $d$):};
  \node[align=left, font=\footnotesize\itshape, prosvBlue] at (-0.5,-2.6) {Проверка (откр.\ ключ $e$):};
\end{tikzpicture}
\caption{Подписание и проверка электронной цифровой подписи на базе RSA.}
\label{fig:p4-2}
\end{figure}

\paragraph*{Проблема большого документа.} В реальности документ может быть огромным (например, несколько мегабайт). Возводить миллион байтов в степень $d$ по модулю $n$ слишком медленно. Поэтому в практических системах поступают иначе:
\begin{enumerate}
  \item вычисляют \term{хеш} (короткий «отпечаток») документа: $h = H(m)$ --- небольшое число, скажем, длиной в 256 бит;
  \item подписывают именно хеш: $\sigma = h^{d} \bmod n$.
\end{enumerate}

При проверке вычисляют $h' = H(m)$ заново и сравнивают $\sigma^{e} \bmod n$ с $h'$. О хеш-функциях и их математическом устройстве пойдёт речь в §\ref{sec:hashing}.

\begin{importantbox}[title={Важно: это учебная схема}]
Здесь показано лишь математическое ядро RSA-подписи. В реальных стандартах подписывают не «голый» хеш, а специально закодированное значение: применяются схемы дополнения (padding) вроде RSASSA-PSS или RSASSA-PKCS1-v1\_5. Без такого кодирования «наивная» подпись $\sigma = h^{d} \bmod n$ уязвима для ряда атак.
\end{importantbox}

\begin{interestbox}[title={Это интересно: ЭЦП в России}]
В Российской Федерации электронная цифровая подпись юридически приравнена к собственноручной с 2002 года (закон №\,1-ФЗ «Об электронной цифровой подписи», заменён в 2011 году на №\,63-ФЗ «Об электронной подписи»). С помощью ЭЦП сегодня сдают налоговую отчётность, регистрируют недвижимость, участвуют в государственных закупках --- всё это без бумажных носителей. Криптографические алгоритмы, используемые в российских ЭЦП, описаны в стандарте ГОСТ Р 34.10--2012; математически они близки к описанной RSA-схеме, но используют не модулярное возведение в степень, а операции на так называемых \emph{эллиптических кривых}.
\end{interestbox}

% ---------------------------------------------------------------------
\subsection{Электронное голосование}
\label{subsec:voting}

С распространением Интернета возникает естественный вопрос: можно ли перевести голосование --- например, выборы или опросы --- в цифровой формат? На первый взгляд, это просто: создать сайт, попросить каждого выбрать вариант, посчитать голоса. На деле возникает несколько противоречивых требований, и решить их одновременно --- содержательная криптографическая задача.

\begin{importantbox}[title={Важно: требования к электронному голосованию}]
Хорошая система электронного голосования должна одновременно обеспечивать:
\begin{enumerate}
  \item \term{Корректность}: каждый имеющий право проголосовать голосует не более одного раза.
  \item \term{Анонимность}: никто, включая организаторов, не может сопоставить конкретного избирателя с его голосом.
  \item \term{Проверяемость}: избиратель может убедиться, что его голос учтён, а наблюдатели --- что общий подсчёт верен.
  \item \term{Неподдельность}: никто не может вбросить «лишние» голоса.
\end{enumerate}
\end{importantbox}

Заметим противоречие между пунктами 1 и 2: чтобы убедиться, что Алиса голосует не дважды, надо как-то её опознать; но если её опознали, как сохранить анонимность? На бумажных выборах эта задача решается \emph{физически} (список явки и тайная урна разделены), а в цифровом мире --- математически. Ключевой инструмент --- \term{слепая цифровая подпись}, предложенная Дэвидом Чаумом ещё в 1982 году.

\paragraph*{Слепая подпись.} Идея слепой подписи: можно подписать документ, \emph{не зная его содержания}. Этот странный, казалось бы, инструмент имеет естественный аналог в офлайне.

\begin{interestbox}[title={Это интересно: аналогия с конвертом и копиркой}]
Представьте: вы кладёте лист бумаги в конверт, внутри которого вложена копирка. Затем просите начальника поставить подпись \emph{на конверте}. Подпись через копирку переходит и на лист внутри. Раскрыв конверт, вы получаете подписанный документ, причём начальник никогда не видел его содержания.

Слепая RSA-подпись --- это математическая реализация той же идеи: «копирка» --- модулярное умножение на случайный множитель.
\end{interestbox}

Формально схема устроена так. Пусть у центра голосования есть RSA-ключи $(e, n)$ (открытый) и $d$ (закрытый).

\begin{algorithmbox}[title={Алгоритм: слепая подпись RSA}]
Алиса хочет получить подпись центра под её бюллетенем $m$, не раскрывая центру содержимое.
\begin{enumerate}
  \item Алиса выбирает случайное $k$, взаимно простое с $n$, и вычисляет \emph{ослеплённое} сообщение
  \[
    m' = m \cdot k^{e} \bmod n.
  \]
  Отправляет $m'$ центру.

  \item Центр (после проверки, что Алиса действительно имеет право голоса) ставит подпись:
  \[
    \sigma' = (m')^{d} \bmod n = m^{d} \cdot k^{ed} \bmod n = m^{d} \cdot k \bmod n.
  \]
  Возвращает $\sigma'$ Алисе.

  \item Алиса \emph{снимает ослепление}: вычисляет
  \[
    \sigma = \sigma' \cdot k^{-1} \bmod n = m^{d} \bmod n.
  \]
  Получена настоящая RSA-подпись центра под её бюллетенем $m$, хотя центр никогда не видел $m$.
\end{enumerate}
\end{algorithmbox}

\paragraph*{Где здесь теория чисел.} Заметим, что на шаге 3 Алиса делит на $k$, точнее --- умножает на $k^{-1} \bmod n$. Существование \emph{обратного} по модулю $n$ обеспечивается тем, что $\gcd(k, n) = 1$, и его можно эффективно вычислить расширенным алгоритмом Евклида (§\ref{sec:numtheory}). Кроме того, в выкладке использовалось $k^{ed} \equiv k \pmod n$ --- следствие малой теоремы Ферма (точнее, её обобщения --- теоремы Эйлера, гласящей, что $k^{\Eul(n)} \equiv 1 \pmod n$ для $\gcd(k, n) = 1$).

\paragraph*{Сам протокол голосования.} На базе слепой подписи строится следующая схема (в упрощённом виде; подробности --- в книге Музыкантского и Фурина).

\begin{algorithmbox}[title={Алгоритм: электронное голосование}]
\begin{enumerate}
  \item \emph{Подготовка.} Алиса формирует бюллетень $m$ --- например, число, кодирующее её выбор, плюс случайный «уникальный идентификатор» $r$, чтобы исключить совпадение бюллетеней.
  \item \emph{Получение подписи.} Алиса ослепляет $m$, отправляет $m'$ \term{Регистратору}. Регистратор проверяет по своему списку, что Алиса имеет право голосовать и ещё не получала подпись, и возвращает $\sigma'$. Алиса снимает ослепление и получает $(m, \sigma)$.
  \item \emph{Анонимная отправка.} Через анонимный канал (например, через сеть посредников или с другого устройства) Алиса отправляет $(m, \sigma)$ \term{Счётчику}.
  \item \emph{Подсчёт.} Счётчик проверяет подпись $\sigma$: если она верна и идентификатор $r$ ранее не встречался --- голос засчитан.
\end{enumerate}
\end{algorithmbox}

\paragraph*{Почему это работает.}
\begin{itemize}
  \item \emph{Корректность.} Регистратор контролирует, что одна и та же Алиса не получит подпись дважды.
  \item \emph{Анонимность.} Регистратор знает, кто такая Алиса, но не видит содержимого $m$ (оно ослеплено). Счётчик видит $m$ и $\sigma$, но не знает, кому принадлежит подпись --- лишь то, что она поставлена настоящим Регистратором.
  \item \emph{Неподдельность.} Без знания закрытого ключа $d$ невозможно сгенерировать $\sigma$ к произвольному $m$ --- это потребовало бы взлома RSA.
  \item \emph{Проверяемость.} Список $\{(m, \sigma)\}$ публикуется. Алиса может убедиться, что её $(m,\sigma)$ в списке; любой наблюдатель может пересчитать сумму, не зная имён.
\end{itemize}

Общую идею протокола поясняет рис.~\ref{fig:p4-3}.

\begin{figure}[H]
\centering
\begin{tikzpicture}[
  >=Stealth, every node/.style={font=\footnotesize},
  voter/.style={draw=prosvBlue, fill=prosvLight, rounded corners, minimum width=2.4cm, minimum height=1cm, align=center, thick},
  server/.style={draw=prosvRed, fill=prosvCream, rounded corners, minimum width=2.4cm, minimum height=1cm, align=center, thick},
]
  \node[voter] (alice) at (0, 0) {Алиса};
  \node[server] (reg) at (5, 1.5) {Регистратор\\ {\scriptsize ключ $d$}};
  \node[server] (count) at (5, -1.5) {Счётчик};

  \draw[->, thick, prosvRed!80] (alice.north east) -- node[above, sloped]{$m' = m \cdot k^e$} (reg.west);
  \draw[<-, thick, prosvRed!80] ([yshift=-15pt]alice.east) -- node[below, sloped]{$\sigma' = (m')^d$} ([yshift=-15pt]reg.west);
  \draw[->, thick, prosvGreen!60!black] (alice.south east) -- node[below, sloped]{анонимно: $(m, \sigma)$} (count.west);
\end{tikzpicture}
\caption{Электронное голосование на базе слепой подписи. С Регистратором Алиса работает «лицом к лицу» (он знает её, но не видит бюллетень). К Счётчику бюллетень приходит анонимно, но с проверяемой подписью Регистратора.}
\label{fig:p4-3}
\end{figure}

\begin{interestbox}[title={Это интересно: голосование в реальности}]
Системы электронного голосования на основе слепой подписи (и более сложных схем) уже применялись на практике: например, в Эстонии часть выборов проходит через интернет с 2005 года, в России на муниципальных выборах в Москве с 2019 года использовалась система на базе криптографических конструкций (правда, на основе не только RSA, но и других схем --- с гомоморфным шифрованием и блокчейном). Подробное обсуждение разных схем электронного голосования можно найти в книге Музыкантского и Фурина «Лекции по криптографии».
\end{interestbox}

% ---------------------------------------------------------------------
\subsection{Итог параграфа}

Мы увидели, что одна и та же конструкция RSA --- модулярное возведение в степень --- решает четыре совершенно разные задачи:
\[
  \boxed{\;\text{шифрование}\;} \qquad
  \boxed{\;\text{аутентификация}\;} \qquad
  \boxed{\;\text{цифровая подпись}\;} \qquad
  \boxed{\;\text{слепая подпись}\;}
\]
Различие лишь в том, \emph{в каком порядке} применяются операции $\cdot^{e}$ и $\cdot^{d}$ и какие именно числа подставляются на вход. В этом одна из главных красот современной криптографии: глубокая математическая идея даёт целый каскад прикладных решений.

\begin{tasksbox}
\begin{enumerate}
  \item Объясните своими словами, в чём принципиальная разница между «шифрованием с открытым ключом» и «цифровой подписью», если в обеих схемах используются одни и те же ключи и одна и та же формула.
  \item Почему в протоколе аутентификации банк каждый раз присылает \emph{новое случайное} $r$, а не одно и то же?
  \item В схеме ЭЦП Алиса публикует пару $(m, \sigma)$. Что произойдёт, если она опубликует пару $(m', \sigma)$ с изменённым $m' \ne m$? Сможет ли она тем самым обмануть проверяющего?
  \item* В схеме слепой подписи покажите подробной выкладкой, что после снятия ослепления получается именно $m^d \bmod n$ (используйте $k^{ed} \equiv k \pmod n$).
  \item* Предложите простую (возможно, наивную) атаку на схему голосования, если из протокола убрать пункт о случайном идентификаторе $r$ в бюллетене.
  \item* Подумайте, какие ещё свойства бумажного бюллетеня (например, возможность переголосовать, защита от принуждения) электронное голосование решает плохо. Каковы возможные подходы?
\end{enumerate}
\end{tasksbox}

% =====================================================================
%   §5. Хеширование, теория чисел и рандомизация
% =====================================================================
\section{Хеширование: теория чисел встречается с рандомизацией}
\label{sec:hashing}

В предыдущих параграфах мы видели, как теория чисел даёт «жёсткие» гарантии безопасности: вычислить нужный ответ практически невозможно из-за вычислительной сложности задачи факторизации. Сейчас мы рассмотрим совсем другой тип использования теории чисел --- \term{вероятностные алгоритмы}, в которых правильность работы достигается не «навсегда», а \emph{в среднем}, с очень высокой вероятностью.

Главная мысль: добавив в алгоритм \emph{случайный выбор} (на этапе настройки, ещё до запуска), можно гарантировать, что для \emph{любых} входных данных алгоритм работает в среднем хорошо. Принципиальное отличие от наивного «алгоритм работает хорошо для случайных данных» в том, что входные данные могут быть какими угодно --- даже подобраны злоумышленником, --- а случайность спрятана \emph{внутри} алгоритма.

Этот параграф следует изложению из классической книги С.\ Дасгупты, Х.\ Пападимитриу и У.\ Вазирани «Алгоритмы»\footnote{Дасгупта~С., Пападимитриу~Х., Вазирани~У. \emph{Алгоритмы}. М.: МЦНМО, 2014. Глава о хешировании.}, в котором роль теории чисел особенно прозрачна.

% ---------------------------------------------------------------------
\subsection{Задача быстрого поиска}

Представьте сервис электронной почты, в котором зарегистрировано сто миллионов пользователей. Каждый пользователь идентифицируется адресом, скажем, длиной до 30 символов. При каждом входе нужно \emph{мгновенно} (за время, не зависящее от размера базы) определить: есть ли такой пользователь, и если есть --- найти его данные.

Технически: имеется большая «вселенная» возможных ключей $U$ (всех допустимых адресов), и интересующее нас подмножество $S \subseteq U$ из $n$ реальных пользователей. Мы хотим организовать структуру данных, позволяющую за $O(1)$ операций проверять, лежит ли $x \in S$, и получать данные, связанные с $x$.

\paragraph*{Прямая адресация.} Самое простое решение --- завести массив, индексируемый \emph{всеми} элементами $U$. Такой массив имеет длину $|U|$. Если $|U|$ велико (а оно \emph{колоссально}: количество всевозможных 30-символьных строк --- это $26^{30}$, что больше $10^{42}$), это решение нереализуемо: не хватит памяти всей планеты.

\paragraph*{Идея хеширования.} Вместо того, чтобы выделять ячейку под каждый возможный элемент, заведём массив размера $m \approx n$ (того же порядка, что число реальных пользователей), и зададим \term{хеш-функцию} $h: U \to \{0, 1, \dots, m-1\}$, которая каждому ключу сопоставляет номер ячейки.

\begin{definition}{хеш-функция и хеш-таблица}{def:p5-1}
\term{Хеш-функция} --- это произвольная функция $h: U \to \{0, 1, \dots, m-1\}$. По ней строится \term{хеш-таблица}: массив $T[0], T[1], \dots, T[m-1]$. Ключ $x \in U$ хранится в ячейке $T[h(x)]$.

Если двум разным ключам $x \ne y$ присвоилось одно и то же значение $h(x) = h(y)$, мы говорим, что произошла \term{коллизия}.
\end{definition}

Иллюстрация работы хеш-функции и появления коллизий приведена на рис.~\ref{fig:p5-1}.

\begin{figure}[H]
\centering
\begin{tikzpicture}[
  >=Stealth, every node/.style={font=\footnotesize},
  uni/.style={draw=prosvBlue, fill=prosvLight, rounded corners, minimum height=0.5cm, minimum width=2.2cm},
  cell/.style={draw=prosvRed!70, fill=prosvCream, minimum height=0.6cm, minimum width=0.7cm},
]
  \node[uni] (k1) at (0, 2)   {ivan@mail};
  \node[uni] (k2) at (0, 1.3) {olga@mail};
  \node[uni] (k3) at (0, 0.6) {petr@yand};
  \node[uni] (k4) at (0, -0.1){anna@bk};
  \node[uni] (k5) at (0, -0.8){vova@mail};

  \node[font=\footnotesize, prosvBlue] at (0, 2.7) {ключи $x$ (вселенная $U$)};

  \foreach \i in {0,1,2,3,4,5,6,7} {
    \node[cell] (c\i) at (5 + 0.75*\i, 0.5) {\i};
  }
  \node[font=\footnotesize, prosvBlue] at (7.5, 1.3) {хеш-таблица $T$ (размер $m=8$)};

  \draw[->, prosvGreen!60!black] (k1.east) -- (c3.west);
  \draw[->, prosvGreen!60!black] (k2.east) -- (c0.west);
  \draw[->, prosvGreen!60!black] (k3.east) -- (c5.west);
  \draw[->, prosvGreen!60!black] (k4.east) -- (c1.west);
  \draw[->, prosvRed!80, very thick] (k5.east) -- (c3.west);

  \node[font=\footnotesize, prosvRed, align=center] at (4, -1.2) {две стрелки в одну\\ ячейку $=$ коллизия};
\end{tikzpicture}
\caption{Хеш-функция «утрамбовывает» большую вселенную ключей в маленькую таблицу. Иногда возникают коллизии.}
\label{fig:p5-1}
\end{figure}

Если коллизий нет (или их мало), мы успешно сжали огромную вселенную ключей в небольшую таблицу. Если же коллизий много, эффективность падает: на одну ячейку приходится много ключей, поиск замедляется.

% ---------------------------------------------------------------------
\subsection{Беда детерминированного хеширования}

Зафиксируем какую-нибудь конкретную хеш-функцию --- например, $h(x) = x \bmod m$ (в предположении, что ключи $x$ представлены как целые числа). На большинстве входов работает прекрасно. Но злоумышленник, зная нашу функцию, может специально подобрать набор ключей $S = \{0, m, 2m, 3m, \dots\}$ --- все они дают одинаковый остаток~$0$. Все они попадут в одну ячейку. Хеш-таблица деградирует до линейного списка, и каждый поиск занимает $O(n)$ вместо $O(1)$.

Этот сценарий не теоретический: реальные сетевые сервисы испытывали атаки именно такого рода --- атаки «алгоритмическим усложнением», когда поток специально подобранных запросов приводил к замедлению сервера в тысячи раз.

\begin{importantbox}[title={Важно: принципиальное наблюдение}]
Для \emph{любой} фиксированной хеш-функции $h$ найдётся «плохой» набор ключей, на котором коллизий очень много. Никакая «гениальная» формула $h$ не спасает.

Спасает другое: \emph{не фиксировать $h$ заранее}, а выбирать её случайно из специального семейства функций.
\end{importantbox}

% ---------------------------------------------------------------------
\subsection{Универсальное семейство хеш-функций}

Идея, восходящая к Картеру и Вегману (1977), состоит в следующем. Договоримся об \emph{целом семействе} хеш-функций $\mathcal{H} = \{h_a\}$. При создании хеш-таблицы выберем из $\mathcal{H}$ \emph{случайную} функцию $h_a$. Эта функция и будет использоваться. Поскольку выбор делается уже после того, как зафиксированы данные (или независимо от них), злоумышленник заранее не знает, какая $h_a$ будет применена, и не может «подстроиться» под неё.

Чтобы это сработало, семейство $\mathcal{H}$ должно обладать специальным свойством.

\begin{definition}{универсальное семейство}{def:p5-2}
Семейство хеш-функций $\mathcal{H} = \{h: U \to \{0, \dots, m-1\}\}$ называется \term{универсальным}, если для любых двух различных ключей $x, y \in U$ ($x \ne y$) вероятность коллизии при случайном выборе $h \in \mathcal{H}$ не превосходит $1/m$:
\[
  \Pr_{h \in \mathcal{H}}[\,h(x) = h(y)\,] \le \frac{1}{m}.
\]
\end{definition}

Поясним: если бы $h$ присваивала каждому ключу совершенно случайный номер ячейки из $m$ возможных, вероятность совпадения двух конкретных ключей в одной ячейке составила бы ровно $1/m$. Универсальное семейство просит лишь того же порядка --- вероятность \emph{не больше} $1/m$.

\begin{theorem}{оценка коллизий}{thm:p5-1}
Если хеш-функция выбрана из универсального семейства, то математическое ожидание числа коллизий заданного ключа $x$ с прочими $n - 1$ ключами таблицы не превосходит $(n-1)/m$. В частности, если $m \ge n$, среднее число коллизий не превосходит $1$.
\end{theorem}

\paragraph*{Доказательство.} Введём индикаторы: $X_y = 1$, если $h(x) = h(y)$, иначе $X_y = 0$. Общее число коллизий ключа $x$ есть $\sum_{y \in S, y \ne x} X_y$. По линейности математического ожидания:
\[
  \mathbb{E}\!\left[\sum X_y\right] = \sum \mathbb{E}[X_y] = \sum \Pr[h(x) = h(y)] \le (n-1) \cdot \frac{1}{m}.
\]
\hfill$\square$

Таким образом, для $m \approx n$ \emph{в среднем} на одну ячейку приходится не более одной--двух коллизий, и время поиска остаётся константным.

\paragraph*{Где встречается случайность.} Подчеркнём: математическое ожидание берётся \emph{по случайному выбору $h$}, а не по случайному выбору ключей! Сами ключи могут быть «худшими из возможных» --- утверждение всё равно справедливо. Это и есть сила рандомизации.

% ---------------------------------------------------------------------
\subsection{Конструкция универсального семейства через теорию чисел}

Теперь самый интересный вопрос: \emph{как же построить} такое семейство $\mathcal{H}$? И вот тут на сцену выходит теория чисел.

\paragraph*{Шаг 1: представление ключа.} Будем считать, что каждый ключ $x \in U$ --- это \emph{вектор} из $k$ небольших чисел: $x = (x_1, x_2, \dots, x_k)$. Например, если ключ --- это IP-адрес вида 192.168.0.1, его естественно представить как четвёрку чисел $(192, 168, 0, 1)$. Если ключ --- это строка, разобьём её на четвёрки или восьмёрки байтов.

Пусть каждое $x_i$ лежит в диапазоне $0 \le x_i < p$, где $p$ --- \term{простое число}, чуть большее размера наших значений.

\paragraph*{Шаг 2: выбор размера таблицы.} Положим $m = p$. (В этом разделе для простоты принимаем такое допущение; в общем случае схема легко обобщается.)

\paragraph*{Шаг 3: семейство $\mathcal{H}$.} Для каждого вектора $a = (a_1, \dots, a_k)$, где $0 \le a_i < p$, определим хеш-функцию:
\[
  h_a(x) = \big( a_1 x_1 + a_2 x_2 + \dots + a_k x_k \big) \bmod p.
\]
Семейство $\mathcal{H} = \{h_a : a \in \{0, \dots, p-1\}^k\}$. Размер семейства --- $p^k$.

\paragraph*{Шаг 4: случайный выбор.} При построении таблицы выбираем $a$ случайно равновероятно из $\{0, \dots, p-1\}^k$ (например, генератором случайных чисел).

\begin{theorem}{Картер--Вегман}{thm:p5-2}
Описанное выше семейство $\mathcal{H}$ \emph{универсально}.
\end{theorem}

\paragraph*{Доказательство.} Возьмём два разных ключа $x = (x_1, \dots, x_k)$ и $y = (y_1, \dots, y_k)$, $x \ne y$. Поскольку $x \ne y$, найдётся хотя бы один индекс $i$, в котором $x_i \ne y_i$; без ограничения общности будем считать $i = 1$.

Коллизия $h_a(x) = h_a(y)$ означает:
\[
  a_1 x_1 + a_2 x_2 + \dots + a_k x_k \equiv a_1 y_1 + a_2 y_2 + \dots + a_k y_k \pmod p,
\]
или, после перестановки,
\[
  a_1 (x_1 - y_1) \equiv -\sum_{i=2}^{k} a_i (x_i - y_i) \pmod p. \qquad (\ast)
\]
Обозначим правую часть как $R$. Заметим, что $R$ зависит от $a_2, \dots, a_k$, но \emph{не зависит} от $a_1$.

\textbf{Здесь и появляется теория чисел.} Поскольку $p$ --- \emph{простое}, а $x_1 - y_1 \ne 0 \pmod p$ (по нашему предположению $x_1 \ne y_1$, и оба они в диапазоне $[0, p)$, значит их разность ненулевая по модулю~$p$), у числа $(x_1 - y_1)$ \emph{существует и единственен} обратный элемент $(x_1 - y_1)^{-1} \bmod p$. Это --- ровно то свойство модулярной арифметики над простым модулем, которое мы доказали в §\ref{sec:numtheory} (расширенный алгоритм Евклида).

Следовательно, уравнение $(\ast)$ имеет ровно одно решение относительно $a_1$ при фиксированных $a_2, \dots, a_k$:
\[
  a_1 = R \cdot (x_1 - y_1)^{-1} \bmod p.
\]
Среди $p$ возможных значений $a_1$ только одно приводит к коллизии. Поэтому условная вероятность коллизии (при фиксированных $a_2, \dots, a_k$) равна $1/p$. Усреднив по $a_2, \dots, a_k$, получаем
\[
  \Pr_a[h_a(x) = h_a(y)] = \frac{1}{p} = \frac{1}{m}.
\]
\hfill$\square$

\begin{importantbox}[title={Важно: почему именно простое $p$}]
Ключевую роль играет то, что $p$ --- \emph{простое}. Будь $p$ составным, у числа $x_1 - y_1$ \emph{могло бы не быть обратного} (если бы оно делилось на общий с $p$ множитель), и тогда уравнение $a_1 \cdot (x_1 - y_1) \equiv R \pmod p$ либо не имело бы решения, либо имело несколько --- и аккуратной оценки $1/p$ не получилось бы.

Так теорема единственности обратного элемента по простому модулю, открытая в античности, оказывается фундаментом современных структур данных, ежесекундно работающих во всех серверах мира.
\end{importantbox}

% ---------------------------------------------------------------------
\subsection{Пример: маленькая хеш-таблица}

Пусть простое $p = 7$ (значит, и $m = 7$), а ключи --- двойки чисел из $\{0, 1, \dots, 6\}$, то есть $k = 2$. Возьмём пять «настоящих» ключей:
\[
  S = \{(1, 2),\ (3, 4),\ (5, 5),\ (2, 6),\ (4, 1)\}.
\]
Случайно выбираем $a = (a_1, a_2)$. Например, пусть случайный генератор выдал $a = (3, 5)$. Тогда $h_a(x) = (3 x_1 + 5 x_2) \bmod 7$:

\begin{center}
\renewcommand{\arraystretch}{1.25}
\begin{tabular}{ccc}
\rowcolor{prosvLight}
ключ $x$ & $3 x_1 + 5 x_2$ & $\bmod 7$ \\
\hline
$(1, 2)$ & $3 + 10 = 13$  & $6$ \\
$(3, 4)$ & $9 + 20 = 29$  & $1$ \\
$(5, 5)$ & $15 + 25 = 40$ & $5$ \\
$(2, 6)$ & $6 + 30 = 36$  & $1$ \\
$(4, 1)$ & $12 + 5 = 17$  & $3$ \\
\end{tabular}
\end{center}

Получаем коллизию: $(3,4)$ и $(2,6)$ попали в ячейку $1$. Если бы случайный генератор выдал другое $a$, например $a = (2, 1)$, картина была бы другой:

\begin{center}
\renewcommand{\arraystretch}{1.25}
\begin{tabular}{ccc}
\rowcolor{prosvLight}
ключ $x$ & $2 x_1 + x_2$ & $\bmod 7$ \\
\hline
$(1, 2)$ & $4$  & $4$ \\
$(3, 4)$ & $10$ & $3$ \\
$(5, 5)$ & $15$ & $1$ \\
$(2, 6)$ & $10$ & $3$ \\
$(4, 1)$ & $9$  & $2$ \\
\end{tabular}
\end{center}

Снова коллизия (правда, в другой ячейке: $(3,4)$ и $(2,6)$ дали $3$). Удивительно: эти два ключа коллидируют при разных $a$. Дело в том, что $(3, 4) - (2, 6) = (1, -2)$, и многие $a$ удовлетворяют $a_1 - 2 a_2 \equiv 0 \pmod 7$. Однако в среднем по всем $a$ вероятность коллизии любых двух конкретных ключей --- ровно $1/7$, как и обещает теорема.

% ---------------------------------------------------------------------
\subsection{Применения хеширования}

Хеширование --- одна из самых востребованных идей в современной информатике. \emph{Универсальное} хеширование особенно важно для хеш-таблиц и быстрого поиска, а \emph{криптографические} хеш-функции --- для проверки целостности, цифровой подписи и блокчейна. Вот несколько мест, где работают эти идеи:

\begin{itemize}
  \item \term{Поиск дубликатов и быстрая выборка.} Все базы данных, все таблицы соответствий в операционных системах, все словари в языках программирования (Python: \texttt{dict}, JavaScript: \texttt{Map}, Java: \texttt{HashMap}) построены на хешировании.
  \item \term{Проверка целостности файлов.} Контрольные суммы при скачивании файла (CRC, MD5, SHA) --- по сути, хеш-функции. Если хеш скачанного файла не совпадает с заявленным, файл повреждён.
  \item \term{Цифровая подпись.} Как мы видели в §\ref{subsec:eds}, RSA-подпись на самом деле ставится не на весь документ, а на его хеш.
  \item \term{Антивирусы и системы обнаружения дубликатов.} База вирусных сигнатур --- это хеши известных вредоносных файлов; чтобы проверить, не заражён ли файл, антивирус считает его хеш и ищет в базе.
  \item \term{Block\-chain.} В криптовалютах хеш каждого блока включается в следующий блок, образуя цепочку, в которой невозможно незаметно подправить «прошлое».
\end{itemize}

\begin{interestbox}[title={Это интересно: криптографические хеш-функции}]
В этом параграфе мы говорили о «простых» хеш-функциях, минимизирующих коллизии. В криптографии используется родственное, но более сильное понятие --- \term{криптографическая хеш-функция}. От неё требуется не только малое число коллизий, но и \emph{вычислительная невозможность найти} два разных входа с одинаковым хешем (\engterm{collision resistance}), а также найти прообраз заданного хеша (\engterm{preimage resistance}). Примеры таких функций --- SHA-256, BLAKE2, в России --- «Стрибог» (ГОСТ Р 34.11--2012). Внутри они устроены сложнее (и не на основе модулярной арифметики), но используют тот же общий принцип: смешать входные биты так, чтобы малейшее изменение во входе приводило к радикальному изменению в выходе.
\end{interestbox}

\begin{tasksbox}
\begin{enumerate}
  \item Для простого $p = 5$, $k = 1$ и $a = 3$ найдите $h_a(x) = a x \bmod p$ для всех $x \in \{0, 1, 2, 3, 4\}$.
  \item Пусть ключи --- пары $(x_1, x_2)$, $p = 11$, $a = (4, 7)$. Найдите коллидирует ли пара $(1, 3)$ с парой $(8, 1)$ при таком $a$.
  \item Объясните, что произойдёт с оценкой $1/m$, если в семействе Картера--Вегмана взять $p$ составным (например, $p = 6$).
  \item Почему в формуле хеш-функции $h_a(x) = (a_1 x_1 + \dots + a_k x_k) \bmod p$ нет «свободного члена» $a_0$, как в линейной функции $a x + b$? Что изменится, если его добавить?
  \item* Покажите, что если выбирать $a$ всего лишь из множества $\{1\}$ (одно значение), то семейство не универсально.
  \item* Покажите, что вероятность того, что у заданного ключа \emph{не будет} коллизий ни с каким из $n-1$ других ключей в таблице размера $m \ge n$, не меньше $1/2$ (используйте неравенство Маркова).
\end{enumerate}
\end{tasksbox}

% =====================================================================
%   §6. Счётчики с короткой памятью HyperLogLog
% =====================================================================
\section{Счётчики с короткой памятью HyperLogLog и эксперимент Стенли Мильграма}
\label{sec:hll}

В этом, заключительном параграфе главы мы рассмотрим ещё один сюжет на стыке теории чисел, теории вероятностей и информатики --- задачу, в которой важно \emph{подсчитывать число различных объектов} в огромном потоке данных, не имея возможности хранить их все. Это пригодится для задачи, на первый взгляд несвязанной с криптографией: измерения «социальных расстояний» в больших сетях, и проверки знаменитой гипотезы Стенли Мильграма о шести рукопожатиях.

Изложение в этом параграфе опирается на главу 7 книги Литвак и Райгородского «Кому нужна математика?»\footnote{Литвак~Н.\,В., Райгородский~А.\,М. \emph{Кому нужна математика? Содержательное введение в математику и computer science}. М.: Манн, Иванов и Фербер, 2017. Гл.~6, 7 и приложения.}, которую мы рекомендуем для самостоятельного изучения.

% ---------------------------------------------------------------------
\subsection{Задача подсчёта различных элементов}

Представьте, что вы работаете в большой интернет-компании и хотите узнать: сколько \emph{уникальных} пользователей зашло сегодня на сайт? Не общее число обращений (его легко посчитать инкрементируя счётчик при каждом запросе), а именно \emph{различных} людей.

\paragraph*{Точное решение.} Простое и очевидное: завести множество (например, хеш-таблицу из §\ref{sec:hashing}), и при каждом приходящем идентификаторе пользователя проверять, есть ли он там; если нет --- добавить и увеличить счётчик. По окончании дня размер множества --- ответ.

\paragraph*{Проблема.} Если число уникальных пользователей --- сотни миллионов или миллиарды, такая таблица занимает гигабайты памяти. Для одной задачи на одном сервере --- допустимо. А если таких задач тысячи, и каждая на каждом из тысяч серверов? Память становится узким местом.

Возникает естественный вопрос: можно ли подсчитать число различных элементов в потоке \emph{приближённо}, но используя \emph{гораздо меньше памяти}? Например, дать ответ с относительной погрешностью $\pm 2\%$, заняв всего несколько килобайт?

Удивительно, но ответ положительный. И построение такого алгоритма опирается на красивые идеи теории вероятностей и --- что нам особенно интересно --- использует хеш-функции, родственные тем, что мы изучали в §\ref{sec:hashing}.

% ---------------------------------------------------------------------
\subsection{Идея: «самый длинный хвост нулей»}

Алгоритм, который мы сейчас опишем (он называется \term{HyperLogLog}, предложен Филиппом Флажоле и его коллегами в 2007 году), основан на одной красивой вероятностной идее.

Пусть у нас есть «хорошая» хеш-функция $h$, которая для каждого входного объекта (например, идентификатора пользователя) выдаёт случайную на вид последовательность из $L$ битов. Слово «случайную» означает: для разных входов биты выглядят как независимые подбрасывания монеты\footnote{Конечно, $h$ детерминированна, но её устройство таково, что заранее, не считая, угадать биты невозможно. На практике используются криптографические или близкие к ним хеш-функции, обладающие именно таким свойством «равномерности».}.

\paragraph*{Наблюдение.} Если число $h(x)$ действительно ведёт себя случайно, то рассмотрим, сколько \emph{нулей подряд} стоит в начале двоичной записи $h(x)$.
\begin{itemize}
  \item Вероятность, что первый бит --- ноль: $1/2$.
  \item Вероятность, что первые два бита --- нули: $1/4$.
  \item Вероятность, что первые $k$ битов --- нули: $1/2^k$.
\end{itemize}

Наглядно эта закономерность показана на рис.~\ref{fig:p6-1}.

\begin{figure}[H]
\centering
\begin{tikzpicture}[every node/.style={font=\small\ttfamily}]
  \node at (0,0)   {010101...};
  \node at (0,-0.5){0000111...};
  \node at (0,-1)  {001011...};
  \node at (0,-1.5){\textcolor{prosvRed}{00000}10...};
  \node at (0,-2)  {1110...};
  \node[font=\small\rmfamily, prosvBlue] at (4, -1.5) {длинный «хвост нулей» --- редкое событие};
  \draw[->, prosvRed] (2.5, -1.5) -- (0.8, -1.5);
\end{tikzpicture}
\caption{Иллюстрация двоичных хешей. Чем длиннее «хвост нулей» в начале, тем реже такое событие.}
\label{fig:p6-1}
\end{figure}

Иначе говоря, если в потоке встречается элемент $x$, такой что $h(x)$ начинается, скажем, с $20$ нулей подряд, это значит, что в потоке \emph{скорее всего побывало порядка $2^{20} \approx 10^6$ различных элементов} --- иначе крайне маловероятно увидеть такое редкое событие.

\paragraph*{Базовый алгоритм («одиночный счётчик»).}

\begin{algorithmbox}[title={Алгоритм: наивный счётчик различных элементов}]
\begin{enumerate}
  \item Заведём одну переменную $M = 0$.
  \item Для каждого приходящего элемента $x$ потока:
  \begin{enumerate}
    \item вычислим $h(x)$;
    \item посчитаем число $\rho(h(x))$ --- длину начального блока нулей в $h(x)$, плюс единица;
    \item обновим $M \leftarrow \max(M, \rho(h(x)))$.
  \end{enumerate}
  \item На выходе оценим число различных элементов как $\widehat{n} = 2^M$.
\end{enumerate}
\end{algorithmbox}

\paragraph*{Память.} Чтобы хранить $M$ от $0$ до $L$ (где $L$ --- длина хеша, скажем, 32 бита), достаточно \emph{шести битов памяти}. Шесть битов --- чтобы оценить число различных элементов в потоке любой длины!

\paragraph*{Проблема.} Эта оценка очень шумная: одно «случайно длинное» обнуление сбивает её в большую сторону, а его отсутствие --- в меньшую. Дисперсия чудовищная.

% ---------------------------------------------------------------------
\subsection{Идея «несколько счётчиков»}

Чтобы сгладить случайные колебания, разделим поток на много подпотоков и усредним результаты. Это и есть алгоритм HyperLogLog.

\begin{algorithmbox}[title={Алгоритм: HyperLogLog (упрощённая схема)}]
\begin{enumerate}
  \item Выберем число $b$ --- скажем, $b = 10$, --- и заведём $2^b = 1024$ независимых счётчика $M_0, M_1, \dots, M_{2^b - 1}$. Каждый счётчик --- 6 битов.
  \item Для каждого приходящего элемента $x$:
  \begin{enumerate}
    \item вычислим $h(x)$ длины $L$;
    \item первые $b$ битов $h(x)$ интерпретируем как номер счётчика $j$, $0 \le j < 2^b$;
    \item к остальным $L - b$ битам применяем правило: $\rho = $ длина начального блока нулей плюс 1;
    \item обновляем: $M_j \leftarrow \max(M_j, \rho)$.
  \end{enumerate}
  \item На выходе оцениваем
  \[
    \widehat{n} = \alpha_b \cdot 2^b \cdot \left( \frac{1}{2^b} \sum_{j=0}^{2^b-1} 2^{-M_j} \right)^{-1},
  \]
  где $\alpha_b$ --- известная константа корректировки (для $b = 10$ примерно $0.7213$).
\end{enumerate}
\end{algorithmbox}

\paragraph*{Почему так?} Формула в третьем шаге --- это \emph{гармоническое среднее} величин $2^{M_j}$, а не простое арифметическое. Гармоническое среднее устойчиво к выбросам (одно случайно большое $M_j$ его существенно не искажает) --- именно поэтому HyperLogLog так точен.

\begin{theorem}{точность HyperLogLog}{thm:p6-1}
Относительная погрешность оценки $\widehat{n}$ по сравнению с истинным значением $n$ имеет порядок $1.04/\sqrt{2^b}$. При $b = 10$ это примерно $3.25\%$; при $b = 14$ --- около $0.81\%$.
\end{theorem}

\paragraph*{Память.} При $b = 14$ имеется $2^{14} = 16\,384$ счётчиков по 6 битов --- всего $\approx 12$ КБ. И этим объёмом мы оцениваем число уникальных элементов в потоке длиной до $\approx 10^{18}$ с точностью около $1\%$! (рис.~\ref{fig:p6-2})

\begin{figure}[H]
\centering
\begin{tikzpicture}[
  every node/.style={font=\footnotesize},
  cell/.style={draw=prosvBlue, fill=prosvLight, minimum width=0.7cm, minimum height=0.7cm},
  hl/.style={draw=prosvRed, fill=prosvCream, minimum width=0.7cm, minimum height=0.7cm, thick}
]
  \node[font=\small, prosvBlue] at (-2, 0.7) {Счётчики:};
  \node[cell] at (-1.3, 0) {3};
  \node[cell] at (-0.6, 0) {1};
  \node[hl]   at (0.1, 0)  {7};
  \node[cell] at (0.8, 0)  {2};
  \node[cell] at (1.5, 0)  {4};
  \node[cell] at (2.2, 0)  {0};
  \node[cell] at (2.9, 0)  {3};
  \node[cell] at (3.6, 0)  {1};
  \node[cell] at (4.3, 0)  {2};
  \node[cell] at (5.0, 0)  {5};
  \node at (5.7, 0)   {$\dots$};

  \node[align=left, font=\footnotesize, prosvRed] at (8, 0) {\textit{редкий «хвост»} ---\\ скорее всего \\ много элементов};
  \draw[->, prosvRed] (8, -0.5) -- (0.4, -0.3);

  \node[font=\footnotesize\itshape, prosvBlue, align=center] at (3, -1.2)
    {усреднение по всем счётчикам с помощью гармонического среднего};
\end{tikzpicture}
\caption{Множество счётчиков HyperLogLog. Каждый запоминает максимальную длину «хвоста нулей» для своего подпотока.}
\label{fig:p6-2}
\end{figure}

\begin{interestbox}[title={Это интересно: где работает HyperLogLog}]
HyperLogLog встроен в Redis (популярную систему кэширования), в Google BigQuery (язык запросов к огромным базам данных), в систему аналитики ВКонтакте, в инфраструктуру PostgreSQL и многих других сервисов. Когда вы видите в личном кабинете рекламной системы оценку «охват: 3{,}5\,млн уникальных пользователей», очень вероятно, что под капотом работает именно этот алгоритм.
\end{interestbox}

% ---------------------------------------------------------------------
\subsection{Эксперимент Стенли Мильграма}
\label{subsec:milgram}

Теперь, когда у нас в руках мощный инструмент для подсчёта уникальных объектов, обратимся к классической задаче из социологии, которую Литвак и Райгородский подробно разбирают в главе 7 своей книги.

\paragraph*{Историческое введение.} В 1960-х годах американский социальный психолог \term{Стенли Мильграм} (1933--1984) задался необычным вопросом: \emph{насколько на самом деле «велик» мир}? Точнее: если выбрать наугад двух человек на разных концах страны, через сколько посредников можно их связать?

В 1967 году Мильграм провёл свой знаменитый эксперимент. Участникам в Канзасе и Небраске были даны письма с просьбой передать их некоему адресату в Бостоне, причём \emph{нельзя} было отправлять напрямую --- только через личного знакомого, который, как кажется участнику, может быть ближе к получателю; тот, в свою очередь, выбирал следующего знакомого, и так далее.

\paragraph*{Удивительный результат.} Большинство писем, дошедших до адресата, преодолевали путь \emph{в среднем за 5--6 пересылок}. Так родилась знаменитая гипотеза «шести рукопожатий»: между любыми двумя людьми на Земле в среднем стоит лишь шесть посредников.

\begin{interestbox}[title={Это интересно: «шесть рукопожатий» в литературе}]
Идея «малого мира» проникла в массовую культуру: пьеса Джона Гуэра «Six Degrees of Separation» (1990), фильм с тем же названием, многочисленные «эстафеты знакомств». В науке термин \term{small-world network} стал важным понятием в теории графов и сетевой социологии.
\end{interestbox}

\paragraph*{Современная проверка.} Эксперимент Мильграма страдал серьёзными методологическими дефектами: 1/\,подавляющее большинство писем (около 80\%) до адресата так и не дошли; 2/\,выборка была очень ограниченной (всего около 300 писем); 3/\,«знакомые» отбирались не случайно, а с предположением о близости к получателю.

В современную эпоху, когда подавляющая часть знакомств и контактов перешла в цифровой формат, задачу можно проверить иначе: \emph{построить граф знакомств целиком} и измерить кратчайшие расстояния. Именно это в 2011 году проделали исследователи в рамках одной известной социальной сети --- они проанализировали данные \emph{всех} её пользователей (на тот момент --- около 720 миллионов человек) и более 69 миллиардов «дружеских связей».

\paragraph*{Что они обнаружили.} Средняя длина кратчайшего пути между двумя случайно выбранными пользователями оказалась равной \emph{4{,}74 рукопожатия}. Иначе говоря, в современных социальных сетях гипотеза Мильграма не просто подтверждается, а даже \emph{улучшается}: между двумя пользователями стоит в среднем менее пяти промежуточных знакомых, а не шесть. Принято округлять это число до 4 --- отсюда выражение «четыре рукопожатия» (рис.~\ref{fig:p6-3}).

\begin{figure}[H]
\centering
\begin{tikzpicture}[
  >=Stealth,
  person/.style={circle, draw=prosvBlue, fill=prosvLight, minimum size=0.7cm, font=\footnotesize, inner sep=0},
]
  \node[person] (a) at (0,0)   {А};
  \node[person] (b) at (1.8,0.8) {?};
  \node[person] (c) at (3.6,0)   {?};
  \node[person] (d) at (5.4,0.8) {?};
  \node[person] (e) at (7.2, 0)  {Я};

  \draw[thick, prosvRed] (a) -- (b) -- (c) -- (d) -- (e);

  \node[font=\footnotesize, prosvBlue, align=center] at (3.6, -1.2)
    {4 рукопожатия = 4 ребра в графе знакомств};
\end{tikzpicture}
\caption{«Четыре рукопожатия»: между Алисой и Яковом стоят всего трое промежуточных знакомых.}
\label{fig:p6-3}
\end{figure}

% ---------------------------------------------------------------------
\subsection{При чём здесь HyperLogLog?}

Возникает вопрос: \emph{как именно} вычислить среднюю длину кратчайшего пути в графе из 720 миллионов вершин и 69 миллиардов рёбер? Прямой подсчёт всех пар (а их $\binom{720{\cdot}10^6}{2} \approx 2{,}6 \cdot 10^{17}$) непосильно даже для самых мощных кластеров.

Здесь и приходит на помощь HyperLogLog (а также его «родственники» --- алгоритмы вероятностного подсчёта). Идея, идущая в реальном анализе, такова:

\begin{algorithmbox}[title={Алгоритм: оценка расстояний в графе через подсчёт «соседей»}]
\begin{enumerate}
  \item Для каждой вершины $v$ заведём HyperLogLog-счётчик $C_v$.
  \item Вначале добавим в $C_v$ только саму вершину $v$. То есть $|C_v| = 1$.
  \item На шаге $t = 1, 2, \dots, T$:
  \begin{itemize}
    \item для каждой вершины $v$ обновим $C_v \leftarrow C_v \cup \bigcup_{u: (u,v) \text{ -- ребро}} C_u^{(t-1)}$,
    \item то есть в $C_v$ добавляются все вершины, до которых дошли соседи $v$ на предыдущем шаге.
  \end{itemize}
  \item Величина $|C_v|$ после $t$ шагов --- это число вершин, до которых от $v$ существует путь длины $\le t$.
\end{enumerate}
\end{algorithmbox}

\paragraph*{Где экономия памяти.} Если бы мы хранили $C_v$ как точное множество, для одной вершины пришлось бы выделить до 720 миллионов битов памяти --- 90 мегабайт. Умножим на 720 миллионов вершин --- получим 65 петабайт. Невозможно.

А поскольку HyperLogLog-счётчик занимает 12 КБ на вершину, всё хранилище умещается в $720{\cdot}10^6 \cdot 12$ КБ $\approx 8{,}6$ ТБ --- объём, который уже умещается в один большой серверный кластер.

\paragraph*{Объединение HyperLogLog-счётчиков.} У алгоритма есть ещё одно прекрасное свойство: \term{HLL-счётчики можно объединять}. Если $C_1$ оценивает мощность множества $A$, а $C_2$ --- мощность $B$, то новый счётчик $C$, у которого $C[j] = \max(C_1[j], C_2[j])$, оценивает мощность $A \cup B$. Это позволяет распределять вычисления по тысячам машин и в конце «склеивать» результаты.

\begin{importantbox}[title={Важно: краткое резюме}]
Современное измерение «социального расстояния» в больших сетях --- это композиция нескольких идей:
\begin{enumerate}
  \item теория графов даёт постановку задачи (кратчайшие расстояния);
  \item теория чисел и хеширование (§\ref{sec:hashing}) обеспечивают «равномерное» распределение элементов;
  \item теория вероятностей и алгоритм HyperLogLog позволяют оценивать мощности множеств за десятки тысяч раз меньшую память;
  \item распределённые вычисления позволяют разнести работу по тысячам компьютеров.
\end{enumerate}
\end{importantbox}

% ---------------------------------------------------------------------
\subsection{От социологии --- к стратегическим выводам}

Эксперимент Мильграма и его цифровое продолжение --- не просто забавный факт. Из него следуют важные вещи:

\begin{itemize}
  \item \emph{Информация распространяется быстро.} Новость, попавшая в социальную сеть, при равномерной передаче «через знакомых» теоретически охватывает всю сеть за 4--5 шагов. На практике это и наблюдается с вирусными мемами.
  \item \emph{Эпидемии распространяются быстро.} В аналогичной модели заражения болезнь, передающаяся при контакте, охватывает большую долю населения за считанные «волны» контактов.
  \item \emph{Цифровая приватность хрупка.} Если ваш друг знает всё о вас, его друг --- через пять рукопожатий --- статистически связан с миллионами людей. Любая утечка приватной информации распространяется по сети быстрее, чем кажется.
  \item \emph{Поиск работы и идей.} Знаменитая работа социолога Марка Грановеттера 1973 года показала, что новые контакты, идеи, рабочие места часто приходят к нам не от близких друзей, а от «слабых связей» --- знакомых второго и третьего уровня. Малый диаметр сети объясняет, почему такая стратегия работает.
\end{itemize}

\begin{interestbox}[title={Это интересно: модель «малого мира» Уоттса--Строгаца}]
Математическая модель социальной сети, объясняющая эффект малого мира, была построена Дунканом Уоттсом и Стивеном Строгацем в 1998 году. Они показали: если у регулярной решётки добавить совсем немного случайных «дальних» рёбер, диаметр графа резко падает с линейного до логарифмического. Социальная сеть устроена именно так: большинство ваших связей --- локальные (одноклассники, коллеги), но достаточно нескольких «дальних» --- например, родственника в другом городе или знакомого по интернету, --- чтобы вся сеть оказалась «компактной».
\end{interestbox}

\begin{tasksbox}
\begin{enumerate}
  \item Объясните, почему в наивном алгоритме одиночного счётчика оценка $\widehat{n} = 2^M$ имеет такую большую дисперсию.
  \item Пусть хеш-функция выдаёт строки по 8 битов. Сколько (в среднем) различных элементов нужно подать на вход, чтобы хотя бы у одного оказалось 6 нулей в начале хеша?
  \item Объясните, почему гармоническое среднее устойчивее к выбросам, чем арифметическое. Приведите пример набора чисел, в котором эти средние существенно различаются.
  \item Оцените, сколько памяти займёт HyperLogLog при $b = 12$, и какова при этом ожидаемая относительная погрешность.
  \item* Почему в алгоритме оценки расстояний в графе нельзя просто хранить $|C_v|$ как число (это же 4 байта)? Зачем хранить целый HLL-счётчик?
  \item* Среднее число рукопожатий в графе оценено как $4{,}74$. Что в этой формулировке играет роль случайной величины, а что --- константы? Какой стандартной операции теории вероятностей соответствует «среднее число»?
\end{enumerate}
\end{tasksbox}

% =====================================================================
%   Подведение итогов главы
% =====================================================================
\section*{Подведение итогов главы}
\addcontentsline{toc}{section}{Подведение итогов главы}
\markboth{Глава \thechapter. Элементы теории чисел и шифрование}{Подведение итогов главы}

\begin{tcolorbox}[
  enhanced, sharp corners, boxrule=0pt,
  colback=prosvCream, leftrule=4pt, colframe=prosvBlue,
  left=12pt, right=12pt, top=10pt, bottom=10pt,
  before skip=0.5em, after skip=1.5em,
  title={\textbf{\color{prosvBlue}\large Главное в этой главе}},
  fonttitle=\bfseries
]
\textbf{Теория чисел.} Простые числа, остатки от деления, сравнения по модулю --- основы модулярной арифметики. Алгоритм Евклида и его расширенная версия позволяют не только находить НОД, но и решать модулярные линейные уравнения. Быстрое возведение в степень за $O(\log k)$ операций позволяет работать с гигантскими числами; та же идея даёт логарифмический алгоритм для чисел Фибоначчи через возведение матрицы $2{\times}2$ в степень. Малая теорема Ферма $a^{p-1} \equiv 1 \pmod p$ --- ключевое утверждение, обеспечивающее корректность RSA. Простых чисел много (закон $\pi(N) \sim N/\ln N$), и проверка простоты быстра (тест Миллера--Рабина, AKS), но \emph{факторизация} составного числа на множители --- задача экспоненциальной сложности.

\textbf{История криптографии.} От шифра Цезаря --- к симметричным шифрам с одним общим ключом --- к проблеме доставки ключа. Идея \term{односторонней функции} (легко в одну сторону, очень сложно в обратную) и \term{функции с секретом} (трапдор) даёт возможность строить асимметричные шифры с открытым и закрытым ключами.

\textbf{RSA и Диффи--Хеллман.} RSA опирается на сложность факторизации $n = pq$: открытый ключ $(n, e)$ позволяет шифровать, закрытый $d = e^{-1} \bmod \Eul(n)$ --- расшифровывать. Корректность доказывается малой теоремой Ферма. Диффи--Хеллман --- протокол выработки общего ключа по открытому каналу: его стойкость основана на сложности задачи дискретного логарифма. Обе схемы уязвимы к атаке «человек посередине» без аутентификации.

\textbf{Применения RSA.} Та же конструкция, применённая в разных порядках, решает четыре задачи: шифрование сообщения, аутентификация пользователя через challenge-response, электронная цифровая подпись (подпись закрытым ключом, проверка открытым) и слепая подпись (основа протоколов электронного голосования).

\textbf{Хеширование.} Универсальное семейство Картера--Вегмана $h_a(x) = (a_1 x_1 + \dots + a_k x_k) \bmod p$ ($p$ --- простое) гарантирует, что вероятность коллизии любых двух ключей при случайном выборе $a$ не превосходит $1/p$. Ключевую роль играет существование и единственность обратного элемента по простому модулю --- результат расширенного алгоритма Евклида.

\textbf{Вероятностные счётчики.} HyperLogLog оценивает число различных элементов в потоке с точностью $\sim 1\%$, используя всего килобайты памяти. В сочетании с теорией графов это позволяет современным сетям подтверждать гипотезу Стенли Мильграма о малом мире: средняя длина «социальной» цепочки в крупных сетях --- около \emph{четырёх рукопожатий}.
\end{tcolorbox}

\subsection*{Что почитать дальше}
\addcontentsline{toc}{subsection}{Что почитать дальше}

\begin{itemize}
  \item \emph{Музыкантский А.\,И., Фурин В.\,В.} Лекции по криптографии. М.: МЦНМО, разные годы. --- Доступное изложение криптографических протоколов, в том числе электронного голосования; рекомендуется к разделу~1.3 и главе~2.
  \item \emph{Литвак Н.\,В., Райгородский А.\,М.} Кому нужна математика? Содержательное введение в математику и computer science. --- Главы 6, 7 и приложения к ним; материал о малых мирах, графах и вероятностных алгоритмах.
  \item \emph{Дасгупта С., Пападимитриу Х., Вазирани У.} Алгоритмы. --- Глава о хешировании, главы об алгоритмах теории чисел.
  \item \emph{Кнут Д.} Искусство программирования, том~2 (Получисленные алгоритмы). --- Классический фундаментальный труд по алгоритмам теории чисел.
  \item \emph{Виноградов И.\,М.} Основы теории чисел. --- Классический российский учебник теории чисел, доступный школьникам старших классов.
\end{itemize}

\subsection*{Большой итоговый проект}
\addcontentsline{toc}{subsection}{Большой итоговый проект}

В качестве длинного исследовательского проекта на каникулы или на год предлагаем выбрать одно из следующих направлений:

\begin{itemize}
  \item \textbf{«Своя» реализация RSA.} На любом удобном языке программирования (Python, Java, C++) реализуйте полный цикл: генерацию пары ключей $(e, n)$ и $d$ (с честным алгоритмом Миллера--Рабина для подбора простых), шифрование, расшифрование, цифровую подпись. Протестируйте на примерах из этой главы.
  \item \textbf{Атака на «слабый» RSA.} Возьмите малое $n$ (50--100 битов) и попытайтесь его факторизовать. Сравните скорость наивного перебора и метода Полларда $\rho$. Подумайте, при каких размерах $n$ ваш ноутбук «сдаётся».
  \item \textbf{Симуляция «малого мира».} Сгенерируйте случайный граф из 1000--10000 вершин по модели Уоттса--Строгаца. Вычислите средние длины кратчайших путей. Сравните с предсказанием теории.
  \item \textbf{Свой HyperLogLog.} Реализуйте упрощённый вариант HyperLogLog ($b = 8$ или $10$). Сгенерируйте поток из миллиона случайных идентификаторов с заранее известным числом уникальных. Оцените точность алгоритма.
  \item \textbf{Стойкость хеш-функции.} Возьмите криптографическую хеш-функцию (например, SHA-1 уже взломан) и реализуйте поиск коллизий простым перебором. Сравните с теоретической оценкой «парадокса дней рождения».
\end{itemize}

